\subsection{Characterization of varieties by their sheaves of functions}

5.4.1. Let X be a topological space, and let Y be a set. Suppose given, for every open set U of X, a set \(\mathcal{L}(\mathrm{U})\) of maps from U to Y. We shall say that the family \(\mathcal{L} = \{\mathcal{L}(\mathrm{U})\}\) is a sheaf of functions with values in Y if it satisfies the following condition:

Let \((\mathbf{U}_{i})_{i\in\mathbf{I}}\) be a family of open subsets of X, with union U, and let f be a mapping from U to Y; for f to belong to \(\mathcal{L}(\mathbf{U})\) it is necessary and sufficient that \(f|\mathbf{U}_{i}\) belongs to \(\mathcal{L}(\mathbf{U}_{i})\) for all i in I.

5.4.2. Let X and Y be two varieties; for every open set \(U \subset X\) let \(\mathcal{L}(U)\) be the set of morphisms from U into Y; then L is a sheaf of functions with values in Y.

When Y = K, we denote by \(C_{X}^{r}\) the sheaf thus defined.

5.4.3. Let X be a topological space and let S be a set of Banach spaces. For every \(E \in S\) let \(F_{E}\) be a sheaf of functions on X with values in E. Suppose that the family of \(F_{E}\) , for \(E \in S\) , satisfies the following condition:

For every \(x \in X\) there exists an open neighborhood U of x, a space

\(E_{0} \in S\) , and a homeomorphism \(\varphi\) of U onto an open subset of \(E_{0}\) , such that, for every open set \(V \subset U\) and every \(E \in S\) , the set \(\mathcal{F}_{E}(V)\) consists of the functions \(g \circ \varphi\) where g ranges over the space \(\mathcal{C}^{r}(\varphi(V); E)\) .

There then exists on X a structure of a manifold of class \(C^{r}\) and of type S, and only one, compatible with the given topology on X and for which \(F_{E}\) is the sheaf of functions of class \(C^{r}\) on X with values in E.

5.4.4. Let X be a topological space and let F be a sheaf of functions with values in K, satisfying the following condition: for every point x of X, there exist an integer n, an open neighborhood U of x, and a homeomorphism \(\varphi\) of U onto an open subset of \(K^{n}\) , such that, for every open set V of U, the set \(\mathcal{F}(V)\) consists of the functions \(g \circ \varphi\) where g ranges over the set of functions of class \(C^{r}\) with values in K on the open set \(\varphi(V)\) of \(K^{n}\) .

There then exists on X a structure of a locally finite-dimensional manifold of class \(C^{r}\) and only one such that \(F = C_{X}^{r}\) .

5.4.5. Let X and X' be two manifolds of class \(C^{r}\), locally finite-dimensional, and let f be a continuous map from X to X'. For f to be a morphism, it is necessary and sufficient that for every open set U' of X' and every function \(g \in \mathcal{C}^{r}(U'; K)\) , the function \(g \circ f\) belongs to \(\mathcal{C}^{r}(U; K)\) , with \(U = f^{-1}(U')\) .

\subsection{Tangent spaces, tangent linear maps}

5.5.1. Let X be a manifold and let \(a\in\mathbb{X}\) . Consider the pairs \((c,h)\) where \(c=(\mathrm{U},\varphi,\mathrm{E})\) is a chart of the manifold X at \(a\) and where \(h\) is an element of E. Two such pairs \((c,h)\) and \((c',h')\) are said to be equivalent if the derivative at \(\varphi(a)\) of the map \(\varphi'\circ\varphi^{-1}\) (which is defined on a neighborhood of \(\varphi(a)\) ) transforms \(h\) into \(h'\) . One thus obtains an equivalence relation between pairs \((c,h)\) and one calls a tangent vector at \(a\) to X a class of equivalent pairs \((c,h)\) (Sets, Chap.~II, § 6, no. 9).

The tangent vectors at \(a\) to X form a set denoted \(\mathrm{T}_{a}(\mathrm{X})\) . If \(c = (\mathrm{U}, \varphi, \mathrm{E})\) is a chart of the manifold X at \(a\) , the map \(\theta_{c}\) from E into \(\mathrm{T}_{a}(\mathrm{X})\) which associates to an element \(h\) of E the tangent vector represented by the pair \((c, h)\) is a bijection. If we transport to \(\mathrm{T}_{a}(\mathrm{X})\) by means of \(\theta_{c}\) the structure of topological K-vector space on E, the structure thus obtained does not depend on the choice of \(c\) and turns \(\mathrm{T}_{a}(\mathrm{X})\) into a Banach space, called the tangent space to X at \(a\) . The dimension (finite or \(+\infty\) ) of \(\mathrm{T}_{a}(\mathrm{X})\) is equal to the dimension of X at \(a\) .

5.5.2. Let X, Y be two varieties, \(f\) be a morphism from X to Y and \(a\) be a point of X. Consider a chart \(c = (U, \varphi, E)\) of X at \(a\) and a chart \(c' = (V, \psi, F)\) of Y at \(b\) with \(f(U) \subset V\) ; the map \(\Phi = \psi \circ f \circ \varphi^{-1}\) of \(\varphi(U)\) to \(F\) is of class \(C^{r}\) , and its derivative \(\mathbf{D}\Phi(\varphi(a))\) at the point \(\varphi(a)\)

is a continuous linear map from E into F. The continuous linear map \(\theta_{c'} \circ D\Phi(\varphi(a)) \circ \theta_{c}^{-1}\) of \(T_{a}(X)\) to \(T_{b}(Y)\) does not depend on the charts \(c\) and \(c'\) chosen; it is denoted \(T_{a}(f)\) and is called the tangent linear map to \(f\) at \(a\) . If \(g\) is a morphism from Y into a manifold Z, we have:

\[
\mathrm{T} _ {a} (g \circ f) = \mathrm{T} _ {f (a)} (g) \circ \mathrm{T} _ {a} (f).
\]

5.5.3. Let \(f:X\to Y\) be a morphism of manifolds. If \(f\) is locally constant, we have \(\mathrm{T}_{a}(f)=0\) for all \(a\) in X. Conversely, if \(\mathrm{T}_{a}(f)=0\) for all \(a\in X\) and if the field K has characteristic 0, then \(f\) is locally constant.

5.5.4. Let U be an open subset of a manifold X and let f be the canonical injection from U into X; for every point a in U, the map \(\mathrm{T}_{a}(f)\) of \(\mathrm{T}_{a}(\mathrm{U})\) to \(\mathrm{T}_{a}(\mathrm{X})\) is an isomorphism of topological vector spaces, by which these two spaces will henceforth be identified.

5.5.5. Let U be an open subset of a Banach space E; if \(\iota\) is the canonical injection of U into E, the triple \(c = (U, \iota, E)\) is a chart of U, and the atlas \(\{c\}\) defines the manifold structure of U. Given a point \(a\) of U, the map \(\theta_c\) is an isomorphism of topological vector spaces from E onto \(T_a(U)\) ; the inverse isomorphism will be denoted \(\zeta_E(a)\) .

Let \(f\) be a map of class \(\mathbf{C}^{r}\) from U into an open subset V of a Banach space F; for every point \(a\) of U, we have:

\[
\mathbf {D} f (a) \circ \zeta_ {\mathbf {E}} (a) = \zeta_ {\mathbf {F}} (f (a)) \circ \mathrm{T} _ {a} (f)
\]

where Df(a) is the derivative of f at a (1.2.1). In other words, the diagram:

\[
\begin{array}{c} \mathrm{T} _ {a} (\mathrm{U}) \xrightarrow {\mathrm{T} _ {a} (f)} \mathrm{T} _ {f (a)} (\mathrm{V}) \\ \Biggl \downarrow \zeta_ {\mathbf {E}} (a) \qquad \qquad \qquad \Biggl \downarrow \zeta_ {\mathbf {F}} (f (a)) \\ \mathrm{E} \xrightarrow {\mathrm{D} f (a)} \mathrm{F} \end{array}
\]

is commutative.

5.5.6. A cotangent vector to X at a, also called a tangent covector or simply a covector at a, is any continuous linear form on the topological vector space \(\mathrm{T}_{a}(\mathrm{X})\) ; these covectors are therefore the elements of the topological dual \(\mathrm{T}_{a}(\mathrm{X})'\) of the tangent space to X at a. This space will be endowed with the topology of uniform convergence on bounded subsets of \(\mathrm{T}_{a}(\mathrm{X})\) ; equipped with this topology, \(\mathrm{T}_{a}(\mathrm{X})'\) is a Banach space called the cotangent space to X at a. It is also denoted by \(\mathrm{T}_{a}^{\prime}(\mathrm{X})\) .

Let \(f\) be a morphism from X into a Banach space E. One calls the differential of \(f\) at \(a\) and denotes by \(d_{a}f\) or \(df_{a}\) the continuous linear map \(\zeta_{\mathbf{E}}(f(a)) \circ \mathbf{T}_{a}(f)\) of \(\mathbf{T}_{a}(\mathbf{X})\) into E. For \(t \in \mathbf{T}_{a}(\mathbf{X})\) , one sometimes denotes \(t(f)\) or t.f the value \(d_{a}f(t)\) of the linear map \(d_{a}f\) at the point t. The map \(f \mapsto d_{a}f\) is linear.

If \(E = K\) , the differential \(d_a f\) of \(f\) at \(a\) is a covector of \(X\) at \(a\) .

Let E, F, G be three Banach spaces and let \((u, v) \mapsto u \cdot v\) be a continuous bilinear map from \(F \times G\) into E. Let f (resp. g) be a morphism from X to F (resp. G). Then the map \(f \cdot g : x \mapsto f(x) \cdot g(x)\) is a morphism from X to E and we have, for every \(t \in \mathrm{T}_{a}(\mathrm{X})\) :

\[
d _ {a} (f. g) (t) = f (a). d _ {a} g (t) + d _ {a} f (t). g (a).
\]

Taking in particular G = E and F = K, the bilinear map in question being multiplication, we have:

\[
d _ {a} (f g) = f (a) d _ {a} g + g (a) d _ {a} f.
\]

5.5.7. Let X be a locally finite-dimensional manifold and let \(\xi^{1},\ldots,\xi^{m}\) be morphic functions defined in an open neighborhood U of a point \(a\) of X. The following conditions are equivalent:

\begin{enumerate}
\def\labelenumi{(\roman{enumi})}
\item
  there exists an open neighborhood V of a contained in U, such that the functions \(\xi^{i}|V\) (for \(1 \leqslant i \leqslant m\) ) form a coordinate system of X in V;
\item
  the differentials \(d_a\xi^i\) for \(1 \leqslant i \leqslant m\) form a basis of \(T_a(X)'\) ;
\item
  the map \(\xi = (\xi^1, \ldots, \xi^m)\) from U into \(K^m\) is étale at \(a\) (see no. 5.7.6).
\end{enumerate}

For the differentials \(d_{a}\xi^{1},\ldots,d_{a}\xi^{m}\) to be linearly independent, it is necessary and sufficient that there exist a neighborhood V of \(a\) , contained in U, such that the functions \(\xi^{i}|V\) are part of a coordinate system of X in V.

5.5.8. Let X be a locally finite-dimensional manifold and let \(\xi = (\xi^1, \ldots, \xi^m)\) be a coordinate system of X in an open set U. Let \(a \in \mathbf{U}\) : we denote by \((\partial_{1,a}, \ldots, \partial_{m,a})\) the basis of the tangent space \(\mathrm{T}_a(\mathrm{X})\) dual to the basis \((d_a\xi^1, \ldots, d_a\xi^m)\) of \(\mathrm{T}_a(\mathrm{X})'\) . We therefore have:

\[
\partial_ {i, a} (\xi^ {j}) = \delta_ {i, j} \quad (\text { Kronecker index }).
\]

The tangent vector \(\partial_{i,a}\) is also denoted \((\partial/\partial\xi^{i})_{a}\) .

Let \(f\) be a function of class \(\mathbf{C}^r\) on U, with values in a Banach space E. We denote by \(\partial_i f\) or \(\partial f / \partial \xi^i\) the function \(a \mapsto \partial_{i,a}(f)\) : this is a function of class \(\mathbf{C}^{r-1}\) on U with values in E (a continuous function if \(r = 1\) ). Its value at a point \(a\) of U is sometimes denoted \((\partial f / \partial \xi^i)_a\) . For every \(t \in \mathrm{T}_a(X)\) , we have:

\[
d _ {a} f (t) = \sum_ {i = 1} ^ {m} d _ {a} \xi^ {i} (t) \frac {\partial f}{\partial \xi^ {i}} (a),
\]

which can also be written:

\[
d _ {a} f = \sum_ {i = 1} ^ {m} \frac {\partial f}{\partial \xi^ {i}} (a) d _ {a} \xi^ {i}.
\]

These notations are consistent with those of 1.6.3.

5.5.9. Let X and Y be two manifolds and let f be a morphism from X to Y. We call the rank of f at a point a of X, and denote by \(rg_{a}f\) , the rank (finite or equal to +∞) of the linear map \(\overline{T}_{a}(f)\) . The map \(a \mapsto rg_{a}f\) is lower semicontinuous.

\subsection{Products of manifolds}

5.6.1. Let X and X' be two sets and let \(c = (\mathrm{U}, \varphi, \mathrm{E})\) (resp. \(c' = (\mathrm{U}', \varphi, \mathrm{E}')\) ) be a chart of X (resp. X'). The triple

\[
(\mathbf {U} \times \mathbf {U} ^ {\prime}, \varphi \times \varphi^ {\prime}, \mathbf {E} \times \mathbf {E} ^ {\prime})
\]

is a chart of the product set \(X \times X'\) , denoted \(c \times c'\) .

5.6.2. Let X and X' be two manifolds of class \(C^{r}\). On the product set X × X' there exists one and only one manifold structure of class \(C^{r}\) such that c × c' is a chart of X × X' for every chart c of X and every chart c' of X'. Equipped with this structure, X × X' is called the product manifold of the manifolds X and X'. The product of any finite number of manifolds is defined similarly. The topology of the manifold X × X' is the product topology of the topologies of X and X'. For a ∈ X and b ∈ X', we have:

\[
\dim_ {(a, b)} (\mathrm{X} \times \mathrm{X} ^ {\prime}) = \dim_ {a} \mathrm{X} + \dim_ {b} \mathrm{X} ^ {\prime}.
\]

5.6.3. Let X and X' be two manifolds, and let X × X' be their product. Let a ∈ X and let a' ∈ X'. The canonical projections

\[
\operatorname{pr} _ {1}: \mathrm{X} \times \mathrm{X} ^ {\prime} \rightarrow \mathrm{X} \quad \text { and } \quad \operatorname{pr} _ {2}: \mathrm{X} \times \mathrm{X} ^ {\prime} \rightarrow \mathrm{X} ^ {\prime}
\]

are morphisms. Let \(\pi_{i}\) ( \(i=1,2\) ) be their tangent linear maps at the point \((a,a')\) . The map

\[
(\pi_ {1}, \pi_ {2}): \mathrm{T} _ {(a, a ^ {\prime})} (\mathrm{X} \times \mathrm{X} ^ {\prime}) \rightarrow \mathrm{T} _ {a} (\mathrm{X}) \times \mathrm{T} _ {a ^ {\prime}} (\mathrm{X} ^ {\prime})
\]

is an isomorphism, which allows one to identify the tangent space to X × X' at (a, a') with the product \(T_{a}(X) \times T_{a'}(X')\).

The injection of \(\mathbf{T}_{a}(\mathbf{X})\) to \(\mathbf{T}_{(a,a^{\prime})}(\mathbf{X}\times \mathbf{X}^{\prime})\) resulting from this identification is the tangent linear map at \(a\) to the morphism \(x\mapsto (x,a^{\prime})\) of \(\mathbf{X}\) to \(\mathbf{X}\times \mathbf{X}^{\prime}\) ; we have an analogous result for the injection of \(\mathbf{T}_{a^{\prime}}(\mathbf{X}^{\prime})\) to \(\mathbf{T}_{(a,a^{\prime})}(\mathbf{X}\times \mathbf{X}^{\prime})\) .

5.6.4. Let W, X and X' be three manifolds, and let \(f: \mathbf{W} \to \mathbf{X}\) , \(f': \mathbf{W} \to \mathbf{X'}\) be two maps. For the map

\[
(f, f ^ {\prime}): \mathbf {W} \rightarrow \mathbf {X} \times \mathbf {X} ^ {\prime}
\]

to be a morphism, it is necessary and sufficient that \(f\) and \(f'\) be morphisms (this justifies the use of the term "product", cf.~Ens., chap.~IV, § 2, no. 4). In this case, if \(a\) is a point of W, we have:

\[
\mathrm{T} _ {a} (f, f ^ {\prime}) = (\mathrm{T} _ {a} (f), \mathrm{T} _ {a} (f ^ {\prime})),
\]

taking into account the identification made above.

5.6.5. Let \(f: \mathbf{X} \to \mathbf{Y}\) and \(f': \mathbf{X}' \to \mathbf{Y}'\) be morphisms of manifolds. Then \(f \times f': \mathbf{X} \times \mathbf{X}' \to \mathbf{Y} \times \mathbf{Y}'\) is a morphism. If \(a \in \mathbf{X}\) and \(a' \in \mathbf{X}'\) , we have:

\[
\begin{array}{l} \mathrm{T} _ {(a, a ^ {\prime})} (f \times f ^ {\prime}) = \mathrm{T} _ {a} (f) \times \mathrm{T} _ {a ^ {\prime}} (f ^ {\prime}) \\ \mathrm{rg} _ {(a, a ^ {\prime})} (f \times f ^ {\prime}) = \mathrm{rg} _ {a} f + \mathrm{rg} _ {a ^ {\prime}} f ^ {\prime}. \end{array}
\]

5.6.6. Let \(X_{1}\) , \(X_{2}\) and Z be three manifolds and f a morphism from \(X_{1} \times X_{2}\) to Z. Let \(a \in X_{1}\) and \(b \in X_{2}\) . The tangent linear map to the partial map \(x \mapsto f(x, b)\) (resp. \(y \mapsto f(a, y)\) ) of \(X_{1}\) (resp. \(X_{2}\) ) into Z is denoted \(T_{(a,b)}^{1}(f)\) (resp. \(T_{(a,b)}^{2}(f)\) ). If one identifies \(T_{(a,b)}(X_{1} \times X_{2})\) with \(T_{a}(X_{1}) \times T_{b}(X_{2})\) , we have:

\[
\mathrm{T} _ {(a, b)} (f). (u, v) = \mathrm{T} _ {(a, b)} ^ {1} (f). u + \mathrm{T} _ {(a, b)} ^ {2} (f). v
\]

for every \(u \in \mathrm{T}_a(\mathrm{X}_1)\) and every \(v \in \mathrm{T}_b(\mathrm{X}_2)\) .

5.6.7. ("Implicit function theorem"). With the hypotheses and notations of the preceding number, suppose moreover that \(\mathrm{T}_{(a,b)}^{2}(f)\) is bijective. Then there exists an open neighborhood U of a in \(X_{1}\) and an open neighborhood V of b in \(X_{2}\) with the following property: for every \(x \in U\) , there exists one and only one point \(g(x)\) of V such that \(f(x, g(x)) = f(a, b)\) , and the map g is a morphism from U to V. For every \(x \in U\) , one has

\[
\mathrm{T} _ {x} (g) = - \left(\mathrm{T} _ {(x, g (x))} ^ {2} (f)\right) ^ {- 1} \circ \mathrm{T} _ {(x, g (x))} ^ {1} (f).
\]

\subsection{Immersions, étale morphisms}

5.7.1. Let X and Y be two manifolds, f a morphism from X to Y and a a point of X. Set \(b = f(a)\) . The following conditions are equivalent:

\begin{enumerate}
\def\labelenumi{(\roman{enumi})}
\item
  The linear map \(T_{a}(f)\) is injective and its image is a closed vector subspace of \(\mathbf{T}_{b}(\mathbf{Y})\) admitting a topological supplement \(^{1}\) ;
\item
  There exist a Banach space F, a closed vector subspace E of F admitting a topological supplement, and charts (U, \(\varphi\) , E) of X at \(a\) and (V, \(\psi\) , F) of Y at \(b\) such that \(f(\mathbf{U}) \subset \mathbf{V}\) and \(\varphi = \psi \circ (f|\mathbf{U})\) .
\item
  There exist an open neighborhood U of a, an open neighborhood V of b containing \(f(\mathbf{U})\) , a manifold Z, a point c of Z, and an isomorphism of manifolds g from \(U \times Z\) onto V such that \(f(x) = g(x, c)\) for all \(x \in U\) .
\item
  There exist an open neighborhood U of a, an open neighborhood V of b and a morphism q from V to X with \(f(\mathbf{U}) \subset \mathbf{V}\) , and \(q(f(x)) = x\) for all \(x \in U\) .
\end{enumerate}

When X and Y are of finite dimension, the preceding conditions are still equivalent to the following:

\begin{enumerate}
\def\labelenumi{(\alph{enumi})}
\setcounter{enumi}{21}
\tightlist
\item
  There exist an open neighborhood U of a, a coordinate system \((\eta^{1},\ldots,\eta^{n})\) of Y in an open neighborhood V of b containing \(f(\mathrm{U})\) and an integer \(m\leqslant n\) such that \(\eta^{j}\circ f=0\) for \(m<j\leqslant n\) and that the functions \(\eta^{1}\circ f,\ldots,\eta^{m}\circ f\) form a coordinate system of X in U.
\end{enumerate}

If properties (i) through (iv) are satisfied, one says that \(f\) is an immersion at \(a\) .

The set of points where \(f\) is an immersion is open in \(X\) ; if this open set is all of \(X\), we say that \(f\) is an immersion.

For \(f\) to be an immersion, it is necessary and sufficient that X can be covered by open sets \(\mathbf{U}_{i}\) such that, for all \(i\), \(f|\mathbf{U}_{i}\) is an isomorphism of \(\mathbf{U}_{i}\) onto a submanifold of Y (cf.~no. 5.8.3).

5.7.2. Examples:

\begin{enumerate}
\def\labelenumi{(\alph{enumi})}
\item
  If X is an open subset of a manifold Y, the canonical injection of X into Y is an immersion.
\item
  Let E and F be two Banach spaces and u a continuous linear map from E into F. Then u is an immersion if and only if u is an isomorphism from E onto a closed vector subspace of F admitting a topological complement.
\end{enumerate}

5.7.3. Let \(f: \mathbf{X} \to \mathbf{Y}\) and \(g: \mathbf{Y} \to \mathbf{Z}\) be two morphisms. If \(f\) and \(g\) are immersions, \(g \circ f\) is an immersion. Conversely, if \(g \circ f\) is an immersion, then \(f\) is an immersion. If \(f: \mathbf{X} \to \mathbf{Y}\) and \(f': \mathbf{X}' \to \mathbf{Y}'\) are immersions, \(f \times f'\) is an immersion.

5.7.4. Suppose that X and Y are finite-dimensional varieties over a field K of characteristic 0. If \(f:X\to Y\) is an injective morphism, the set of points of X where \(f\) is an immersion is a dense open subset of X. If K = R or C, this result remains true if one assumes only that X is finite-dimensional.

5.7.5. Let \(f:X\to Y\) be an immersion and let \(g\) be a continuous map from a manifold \(Z\) to \(X\) . For \(g\) to be a morphism, it is necessary and sufficient that \(f\circ g\) be a morphism.

5.7.6. Let \(f: X \to Y\) be a morphism, and let \(a\) be in \(X\) . The following two properties are equivalent:

\begin{enumerate}
\def\labelenumi{(\roman{enumi})}
\item
  \(\mathbf{T}_{a}(f)\) is bijective.
\item
  There exists an open neighborhood U of a and an open neighborhood V of \(f(a)\) such that f induces an isomorphism from U onto V.
\end{enumerate}

When these properties hold, we say that \(f\) is a local isomorphism at \(a\) , or that \(f\) is étale at \(a\) . If this holds for every \(a \in X\) , one says that \(f\) is étale, or is an étale map, or that \(X\) is étale over \(Y\) (by means of \(f\) ). For a morphism to be étale, it is necessary and sufficient that it be both an immersion ( \(n^{\circ}\) 5.7.1) and a submersion ( \(n^{\circ}\) 5.9.1.).

5.7.7. Let \(f:X\to Y\) be an immersion; suppose the variety X is pure of finite dimension. Then \(f\) is étale in each of the following two cases:

\begin{enumerate}
\def\labelenumi{(\roman{enumi})}
\item
  \(\dim Y = \dim X\) ;
\item
  \(f\) is surjective and the topology of X has a countable base of open sets.
\end{enumerate}

5.7.8. For a morphism \(f\) to be an isomorphism onto an open submanifold (resp. an isomorphism), it is necessary and sufficient that \(f\) be étale and injective (resp. étale and bijective).

\subsection{Inverse images of manifold structures, submanifolds}

5.8.1. Let X be a topological space, Y be a manifold, and f be a map from X to Y. Consider the following conditions:

(QR) (resp. (R)) For every \(a \in X\) , there exists an open neighborhood \(U\) of a in \(X\) and a chart \((V, \varphi, E)\) of the manifold \(Y\) at \(f(a)\) such that \(f(U) \subset V\) and that \(\varphi \circ f\) induces a homeomorphism of \(U\) onto the intersection of \(\varphi(V)\) with a closed vector subspace (resp. a closed vector subspace admitting a topological complement) \(F\) of \(E\) .

If condition (QR) is satisfied, there exists on X one and only one manifold structure for which the triples (U, \(\varphi \circ f\) , F) (with the notations of (QR)) are charts of X. It is called the inverse image structure of the manifold structure of Y by \(f\) . Its topology is that of X.

For there to exist on X a manifold structure compatible with the given topology and for which \(f\) is an immersion, it is necessary and sufficient that condition (R) be satisfied. This structure is then unique: it is the inverse image under f of the manifold structure of Y.

5.8.2. The condition (R) above is satisfied in particular when for every \(a \in X\) there exists an open neighborhood U of \(a\) such that \(f|U\) is a homeomorphism from \(U\) onto an open subset of \(Y\). In this case, the manifold \(X\) obtained is étale over \(Y\) (5.7.6).

5.8.3. Suppose now that X is a topological subspace of Y, with f being the canonical injection. If f satisfies condition (R) (resp. (QR)) of no. 5.8.1, then X, equipped with the inverse image structure of the manifold structure of Y under f, is said to be a submanifold (resp. quasi-submanifold) of Y. A submanifold is a quasi-submanifold.

Every quasi-submanifold is locally closed; every open set is a submanifold and its structure is that defined in No. 5.2.3.

When Y is locally of finite dimension, a topological subspace X of Y is a submanifold of Y if and only if, for every \(a \in X\) there exists an open neighborhood U of a in Y, a coordinate system \((\xi^{1}, \ldots, \xi^{m})\) of Y in U, and an integer \(k \leqslant m\) such that \(U \cap X\) is the set of points of U where the functions \(\xi^{i}\) for \(1 \leqslant i \leqslant k\) vanish simultaneously.

Let \(f:X\to Y\) be a morphism of varieties. Suppose that \(f\) is an immersion, and let \(f\) induce a homeomorphism from X onto \(f(X)\) . Then \(f(X)\) is a submanifold of Y, and \(f\) induces an isomorphism from X onto \(f(X)\) ; one then says that \(f\) is an embedding of X into Y.

5.8.4. Let Y and W be two varieties, let X be a subvariety of Y, and let f be a map from X to W. For f to be a morphism, it is necessary and sufficient that, for every \(a \in X\) , there exist a neighbourhood U of a in Y and a morphism \(g : U \to W\) such that g coincides with f on \(U \cap X\) .

5.8.5. Let X be a quasi-subvariety of a variety Y, and let g be a map from a variety Z into X. Suppose that K is of characteristic 0, or that X is a subvariety of Y. For g to be a morphism from Z to X, it is necessary and sufficient that it be a morphism from Z to Y: in other words, the variety structure of X is induced by that of Y (Ens., ch.~IV, § 2, no. 4).

5.8.6. If X is a subvariety (resp. quasi-subvariety) of Y and if Y is a subvariety of a variety Z, then X is a subvariety (resp. quasi-subvariety) of Z.

5.8.7. Let \(X_{i}\) (for i = 1, 2) be a variety and let \(Y_{i}\) be a subvariety (resp. quasi-subvariety) of \(X_{i}\) (for i = 1, 2). Then \(Y_{1} \times Y_{2}\) is a subvariety (resp. quasi-subvariety) of \(X_{1} \times X_{2}\) .

5.8.8. Let X be a quasi-subvariety of a variety Y; let x be a point of X and denote by f the canonical injection of X into Y. The map \(T_{x}(f):T_{x}(X)\to T_{x}(Y)\) is injective and allows one to identify the space \(T_{x}(X)\) with a closed vector subspace of \(T_{x}(Y)\) . The quotient topological vector space \(T_{x}(Y)/T_{x}(X)\) is a Banach space, which is called the transversal space to X at x (in Y). Its dimension (finite or +∞) is called the codimension of X in Y at the point x.

If moreover X is a subvariety of Y, then the space \(T_{x}(X)\) admits a topological supplement in \(T_{x}(Y)\) .

5.8.9. Let \(f\) be a map from a variety X to a variety Y, and let \(\Gamma\) be its graph. For \(f\) to be a morphism, it is necessary and sufficient that the following two conditions be satisfied:

\begin{enumerate}
\def\labelenumi{(\roman{enumi})}
\item
  \(\Gamma\) is a subvariety of \(X \times Y\) .
\item
  For every \((x, y) \in \Gamma\) , one has
\end{enumerate}

\[
\mathrm{T} _ {(x, y)} (\mathrm{X} \times \mathrm{Y}) = \mathrm{T} _ {(x, y)} (\Gamma) \oplus \mathrm{T} _ {y} (\mathrm{Y}).
\]

If this is so, the map \(\mathsf{pr}_1\) induces an isomorphism from \(\Gamma\) onto \(X\) , and \(\mathrm{T}_{(x,y)}(\Gamma)\) is identified with the graph of \(\mathrm{T}_x(f)\) .

In particular, the diagonal of \(X \times X\) is a subvariety of \(X \times X\) .

5.8.10. Let Y be a variety, and let \((f_{i})_{i\in I}\) be a finite family of morphic functions on Y. Let X be the set of \(x\in Y\) such that \(f_{i}(x)=0\) for every i. Assume the following hypothesis:

\begin{enumerate}
\def\labelenumi{(\Alph{enumi})}
\setcounter{enumi}{9}
\tightlist
\item
  For every \(x \in X\) , the differentials \(d_x f_i\) are linearly independent in \(T_x'(Y)\) .
\end{enumerate}

Then X is a closed subvariety of Y, the tangent space \(T_{x}(X)\) is the subspace of \(T_{x}(Y)\) formed by the \(\alpha\) such that \(\alpha.f_{i}=0\) for every i in I. Moreover, the codimension of X in Y is equal to Card (I) at each of its points.

5.8.11. (Simple points of an ideal). Let \(\mathfrak{a}\) be an ideal of the polynomial algebra \(\mathrm{K}[\mathrm{X}_1,\ldots ,\mathrm{X}_n]\) . A point \(x = (x_{1},\dots,x_{n})\) of \(\mathbf{K}^n\) is called a zero of \(\mathfrak{a}\) if \(f(x) = 0\) for all \(f\in \mathfrak{a}\) . If \(x\in \mathbf{K}^n\) , denote by \(\mathbf{S}_x\) the subalgebra of \(\mathrm{K}(\mathrm{X}_1,\dots,\mathrm{X}_n)\) formed by the fractions \(f / g\) , with \(f,g\in \mathrm{K}[\mathrm{X}_1,\dots,\mathrm{X}_n]\) , and \(g(x)\neq 0\) ; denote by \(\mathfrak{a}_x\) the ideal of \(\mathbf{S}_x\) generated by \(\mathfrak{a}\) in \(\mathbf{S}_x\) . A point \(x\) is called a simple zero of \(\mathfrak{a}\) if it is a zero of \(\mathfrak{a}\) and if the following condition is satisfied:

\begin{enumerate}
\def\labelenumi{(\Alph{enumi})}
\setcounter{enumi}{18}
\tightlist
\item
  There exists a finite sequence \((f_1, \ldots, f_m)\) of elements of \(\mathfrak{a}\) which generate the ideal \(\mathfrak{a}_x\) and whose differentials at \(x\) are linearly independent. (This condition is equivalent to saying that the local ring \(\mathbf{S}_x / \mathfrak{a}_x\) is regular (Alg. Comm., to appear).)
\end{enumerate}

Let \(Z\) (resp. \(Z_s\) ) be the set of zeros (resp. simple zeros) of \(\mathfrak{a}\). The set \(Z\) is closed in \(K^n\) , \(Z_s\) is open in \(Z\) , and \(Z_s\) is a subvariety of \(K^{n}\) . If \(x \in Z_{s}\) , the ideal \(\mathrm{K}[\mathrm{X}_{1}, \ldots, \mathrm{X}_{n}] \cap \mathfrak{a}_{x}\) consists of the polynomials f vanishing on a neighborhood of x in \(Z_{s}\) .

Let \(\overline{\mathfrak{a}}\) be the ideal of polynomials vanishing on Z. If K is algebraically closed, the set of simple zeros of \(\mathfrak{a}\) is dense in Z.

5.8.12. Let X be a variety and L be the set of pairs \((x, Z)\) where x is a point of X and Z a subvariety of X containing x. Given two pairs \(\pi = (x, Z)\) and \(\pi' = (x', Z')\) , we denote by \(R\{\pi, \pi'\}\) the relation:

"We have \(x = x'\) and there exists a neighborhood U of \(x\) such that \(\mathbf{U} \cap \mathbf{Z} = \mathbf{U} \cap \mathbf{Z'}\) . Then R is an equivalence relation on L; we denote by \(\gamma_x(Z)\) the equivalence class of the pair \((x, Z) \in L\) . On the set \(J = L/R\) , there exists one and only one variety structure satisfying the following condition:

For every subvariety Z of X, the map \(x \mapsto \gamma_{x}(Z)\) from Z into J is an isomorphism of Z onto an open subvariety of J.

We say that J is the variety of germs of subvarieties of X (cf.~Top. Gén., chap.~I, 4 \(^{e}\) ed., § 6, no. 10).

The map \(\rho: J \to X\) defined by \(\rho(\gamma_x(Z)) = x\) is an immersion; it is called the canonical immersion of \(J\) to \(X\) .

If X is a separated analytic variety of finite dimension, so is J.

\subsection{Submersions and quotient varieties}

5.9.1. Let \(f: X \to Y\) be a morphism of varieties, \(a\) a point of \(X\) and let us set \(b = f(a)\) . The following conditions are equivalent:

\begin{enumerate}
\def\labelenumi{(\roman{enumi})}
\item
  The linear map \(T_{a}(f)\) is surjective, and its kernel admits a topological complement in \(T_{a}(X)\) .
\item
  There exists a chart (U, \(\varphi\) , E) of X at \(a\) , a chart (V, \(\psi\) , F) of Y at \(b\) and a surjective continuous linear map \(u\) from E to F such that
\end{enumerate}

\[
f (\mathrm{U}) \subset \mathrm{V}, \quad \psi \circ f = u \circ \varphi
\]

and that the kernel of \(u\) admits a topological supplement in E.

\begin{enumerate}
\def\labelenumi{(\roman{enumi})}
\setcounter{enumi}{2}
\item
  There exist an open neighborhood U of a, an open neighborhood V of b containing \(f(\mathbf{U})\) and a morphism g from U into a variety Z such that the map \((f,g)\) from U into \(V \times Z\) is an isomorphism.
\item
  There exists an open neighborhood V of b and a morphism s from V to X such that \(s(b) = a\) and \(f(s(y)) = y\) for all y in V (“local section”).
\end{enumerate}

When X and Y are finite-dimensional, the preceding conditions are equivalent to the following condition:

\begin{enumerate}
\def\labelenumi{(\alph{enumi})}
\setcounter{enumi}{21}
\tightlist
\item
  There exist an open neighborhood U of a, an open neighborhood V of b containing \(f(\mathbf{U})\) , and a coordinate system \((\eta^{1},\ldots,\eta^{n})\) on V such that the functions \(\eta^{i}\circ f\) on U be part of a coordinate system on U.
\end{enumerate}

If properties (i) through (iv) are satisfied, one says that \(f\) is a submersion at \(a\) . The set of points where \(f\) is a submersion is open in \(X\) ; if this open set is all of X, then f is said to be a submersion.

5.9.2. Let \(f: X \to Y\) and \(g: Y \to Z\) be two morphisms. If \(f\) and \(g\) are submersions, the same holds for \(g \circ f\) . Conversely, if \(g \circ f\) is a submersion and if \(f\) is surjective, then \(g\) is a submersion.

5.9.3. If \(f: X \to Y\) and \(f': X' \to Y'\) are submersions, \(f \times f'\) is a submersion.

5.9.4. A submersion \(f: X \to Y\) is an open mapping (cf. Top. Gen., chap. I, 4 \(^{e}\) ed., § 5, no. 1); in particular, the equivalence relation R defined by \(f\) is open, \(f\) defines, by passing to the quotient, a homeomorphism from X/R onto \(f(X)\) , and \(f(X)\) is open in Y.

5.9.5. Let R be an equivalence relation on a variety X. We say that R is regular if there exists on the quotient space X/R a variety structure such that the canonical projection \(p: X \to X/R\) is a submersion; this variety structure is then unique; it is a quotient of that of X (Sets, ch.~IV, § 2, no. 6): in other words, for a morphism \(g\) from X/R into a variety Y to exist, it is necessary and sufficient that \(g \circ p\) be a morphism from X to Y.

Let \(C \subset X \times X\) be the graph of R. For R to be regular, it is necessary and sufficient that the following two conditions be satisfied:

\begin{enumerate}
\def\labelenumi{(\roman{enumi})}
\item
  C is a submanifold of X × X.
\item
  The map \(pr_{1}\) from C into X is a submersion.
\end{enumerate}

(ii) also means that if \(a\) and \(b\) are congruent modulo \(\mathbf{R}\) there exists an open neighborhood U of \(a\) and a morphism \(s\) from U into X such that \(s(a) = b\) and \(s(x)\) is congruent to \(x\) modulo \(\mathbf{R}\) for all \(x \in \mathbf{U}\) .

Suppose R is regular. For the quotient variety X/R to be separated, it is necessary and sufficient that the graph of R be closed in \(X \times X\) .

5.9.6. Let X and X' be two varieties, R and R' regular equivalence relations on X and X', and let \(f: X \to X'\) be a morphism compatible with the relations R and R'. The map \(\bar{f}: X/R \to X'/R'\) deduced from \(f\) by passing to quotients is then a morphism.

5.9.7. (“Transitivity of quotients”) Let R and S be two equivalence relations on a variety X such that R implies S, and let S/R be the quotient equivalence relation on X/R. Suppose that R is regular. Then, for S to be regular, it is necessary and sufficient that S/R be regular; if this is the case, the canonical bijection

\[
(\mathrm{X} / \mathrm{R}) / (\mathrm{S} / \mathrm{R}) \rightarrow \mathrm{X} / \mathrm{S}
\]

is an isomorphism of varieties.

5.9.8. ("Products of quotients") Let \((\mathbf{X}_{i})_{i\in I}\) be a finite family of varieties, each equipped with a regular equivalence relation \(R_{i}\) . Let \(X = \prod_{i\in I} X_{i}\) and let R be the equivalence relation on X that is the product of the \(R_{i}\) (cf. Top. Gen., chap. I, 4th ed., § 5, no. 3, cor. of prop. 8). Then R is regular, and the canonical bijection from X/R onto \(\prod_{i\in I}(X_{i}/R_{i})\) is an isomorphism of varieties.

\subsection{Subimmersions}

5.10.1. Let \(f:X\to Y\) be a morphism of manifolds and let \(\Gamma\) be the graph of \(f\) . The map \(j:x\mapsto(x,f(x))\) is an immersion of \(X\) into \(X\times Y\) whose image is the submanifold \(\Gamma\) , and \(f=\mathrm{pr}_{2}\circ j\) is the composite of the immersion \(j\) followed by the submersion \(\mathrm{pr}_{2}\) .

Let \(a \in X\) . We say that \(f\) is a subimmersion at \(a\) if there exists an open neighborhood \(U\) of \(a\) , a manifold \(Z\) , a submersion \(s\) of \(U\) into \(Z\) and an immersion \(i\) of \(Z\) into \(Y\) such that \(f|U = i \circ s\) . The set of points of \(X\) where \(f\) is a subimmersion is an open subset of \(X\) ; if this open set is \(X\) itself, we say that \(f\) is a subimmersion.

5.10.2. Immersions and submersions are subimmersions. If \(f: X \to Y\) is a subimmersion, \(g: Y \to Z\) an immersion, and \(h: W \to X\) a submersion, then \(g \circ f \circ h\) is a subimmersion.

If \(f\) and \(f'\) are subimmersions, \(f \times f'\) is a subimmersion.

5.10.3. For \(f: X \to Y\) to be a subimmersion at a point \(a\) of \(X\) , it is necessary and sufficient that there exist a chart \((U, \varphi, E)\) of \(X\) at \(a\) , a chart \((V, \psi, F)\) of \(Y\) at the point \(f(a)\) , and a continuous linear map \(g\) of \(E\) to \(F\) such that:

\begin{enumerate}
\def\labelenumi{(\roman{enumi})}
\item
  \(f(\mathbf{U}) \subset \mathbf{V}\) , \(g(\varphi(\mathbf{U})) \subset \psi(\mathbf{V})\) and \(f|\mathbf{U} = \psi^{-1} \circ g \circ \varphi\) ;
\item
  the kernel (resp. image) of g is a closed subspace of E (resp. F) admitting a topological supplement.
\end{enumerate}

5.10.4. Let X and Y be two finite-dimensional manifolds. For a morphism \(f\) from X to Y to be a subimmersion at a point \(a\) of X, it is necessary and sufficient that there exist a coordinate system \((\xi^{1},\ldots,\xi^{m})\) of X at \(a\) , a coordinate system \((\eta^{1},\ldots,\eta^{n})\) of Y at \(f(a)\) and an integer \(k\leqslant\inf(m,n)\) such that

\[
\eta^ {i} \circ f = \xi^ {i} \quad \text { for } 1 \leqslant i \leqslant k
\]

\[
\eta^ {i} \circ f = 0 \quad \text {   for   } k <   i \leqslant n.
\]

5.10.5. Let \(f: X \to Y\) be a subimmersion. For every \(b \in Y\) , \(f^{-1}(b)\) is a subvariety of \(X\) ; the subspace of \(\mathrm{T}_{a}(X)\) tangent to the submanifold \(f^{-1}(b)\) at the point \(a \in f^{-1}(b)\) is the kernel of \(\mathrm{T}_{a}(f)\) .

5.10.6. ("Constant rank theorem") Let \(f: X \to Y\) be a morphism of manifolds and let \(a \in X\) . If \(f\) is a subimmersion at \(a\) , one has \(\operatorname{rg}_{x} f = \operatorname{rg}_{a} f\) for \(x\) near \(a\) .

Conversely, suppose the field K has characteristic zero. Let (U, \(\varphi\) , E) be a chart of X at \(a\) and let (V, \(\psi\) , F) be a chart of Y at \(f(a)\) , with \(f(U) \subset V\) . Set \(g = \psi \circ f \circ \varphi^{-1}\) . If there exist a closed vector subspace \(E_{1}\) of E and a closed vector subspace \(F_{1}\) of F such that for every \(x \in U\) , the subspace \(E_{1}\) (resp. \(F_{1}\) ) is a topological supplement of the kernel (resp. image) of the derivative Dg( \(\varphi(x)\) ) of \(g\) at the point \(\varphi(x)\) , then \(f\) is a subimmersion at \(a\) .

If K has characteristic zero (resp. K = R or C) and if Y is finite-dimensional (resp. \(rg_{a}f < +\infty\) ), then f is a subimmersion at a if and only if \(rg_{x}f = rg_{a}f\) for every x sufficiently close to a.

If K has characteristic zero (resp. K = R or C) and Y (resp. X) is finite-dimensional, the set of points \(x \in X\) where f is a subimmersion is a dense open subset of X.

5.10.7. ("Canonical factorization of a subimmersion") Every subimmersion is the composite of a submersion and an immersion. Precisely, let \(f: X \to Y\) be a subimmersion, and let J be the manifold of germs of submanifolds of Y (5.8.12). Let \(x\) be in X and \(y = f(x)\) ; there exist open neighborhoods U of \(x\) such that \(f|U\) is a submersion of U onto a submanifold Z of Y; the element \(\gamma_y(Z)\) of J depends only on \(x\) and not on the chosen neighborhood U, and if we denote it \(\lambda(x)\) , the map \(\lambda\) is a submersion from X into J; if we denote \(\rho\) the canonical immersion of J into Y, we have \(f = \rho \circ \lambda\) . If \(f\) is an immersion, the morphism \(\lambda\) from X into J is étale.

If g is a surjective submersion from X onto a variety Z and h is a morphism from Z to Y such that \(f = h \circ g\) , there exists a unique submersion \(\mu\) from Z into J such that \(\lambda = \mu \circ g\) .

\subsection{Fibered products and inverse images}

5.11.1. Let F be a Banach space, \((\mathrm{E}_{i})_{i\in I}\) a finite family of Banach spaces and \(\mathbf{f}=(f_{i})_{i\in I}\) a family of continuous linear maps \(f_{i}\) of \(E_{i}\) into F. Let E be the product of the \(E_{i}\) and let f be the continuous linear map from E into \(F^{I}\) which is the product of the \(f_{i}\). Finally, let D be the closed subspace of \(F^{I}\) formed by the \((y_{i})_{i\in I}\) such that \(y_{i}\) is independent of i (the “diagonal” of \(F^{I}\) ). The family f is said to be transversal if the continuous linear map obtained by composing f with the canonical projection of \(F^{I}\) onto the quotient \(F^{I}/D\) is surjective and if its kernel \(f^{-1}(D)\) admits a topological complement.

If the spaces \(E_{i}\) and F are all finite-dimensional, the family f is

transversal if and only if one has:

\[
\operatorname{codim} \left(\bigcap_ {i} \operatorname{Im} f _ {i}\right) = \sum_ {i} \operatorname{codim} \left(\operatorname{Im} f _ {i}\right)
\]

If \(I = \{1, 2\}\) , the pair \((f_1, f_2)\) is transversal if and only if the map

\[
f _ {1} + f _ {2}: \mathbf {E} _ {1} \oplus \mathbf {E} _ {2} \rightarrow \mathbf {F}
\]

is surjective and if its kernel admits a topological supplement (if \(E_{1}\) and \(E_{2}\) are of finite dimension, it is equivalent to say that

\[
\operatorname{Im} f _ {1} + \operatorname{Im} f _ {2} = \mathrm{F}).
\]

5.11.2. Let S be a variety, \((X_{i})_{i\in I}\) a finite family of varieties and \(\mathbf{f}=(f_{i})_{i\in I}\) a family of morphisms \(f_{i}\) of \(X_{i}\) into S. Let P be the subset of the product X of the \(X_{i}\) consisting of the points \((x_{i})_{i\in I}\) such that \(f_{i}(x_{i})\) is independent of i. Let \(x\in P\) and let \(y=f_{i}(x_{i})\in S\) . One says that the family f is transversal at the point x of P if the maps \(T_{x}(f_{i})\) form a transversal family of continuous linear maps with values in the Banach space \(T_{y}(S)\) . We say that f is transversal if it is transversal at every point of P.

If f is transversal, then P is a submanifold of X, called the fibered product of the family of \(X_{i}\) above S (relative to the family f), which we denote by \(\prod_{i\in I}X_{i}\) (or more simply \(X_{1}\times_{S}X_{2}\) when \(I=\{1,2\}\) for example). For every point \(x=(x_{i})\) of P, the tangent space \(T_{x}(\prod_{S}X_{i})\) is the subspace of \(\prod T_{x}(X_{i})\) consisting of tangent vectors \(t=(t_{i})\) such that \(T_{x_{i}}(f_{i}).t_{i}\) is independent of the index i.

If \(f_{1}: X_{1} \to Y\) is a submersion and \(f_{2}: X_{2} \to Y\) a morphism, the pair \((f_{1}, f_{2})\) is transversal.

5.11.3. ("Universal property of fibered products") Let \(\mathbf{f} = (f_i)_{i \in I}\) be a transversal family of morphisms \(f_i: X_i \to S\) and let P be the fibered product of the \(X_i\) over S; for every \(i \in I\) , one denotes by \(\pi_i\) the morphism from P into \(X_i\) obtained by restricting to P the projection of X onto \(X_i\) ; then \(f_i \circ \pi_i\) is a morphism from P to S independent of \(i\) . Let T be a variety, and let \(g_i: T \to X_i\) be morphisms of varieties such that \(f_i \circ g_i\) is a morphism from T to S independent of \(i \in I\) ; there then exists a unique morphism \(h\) from T to P such that \(g_i = \pi_i \circ h\) for all \(i \in I\) .

5.11.4. ("Associativity of the fibered product") Let \(\mathbf{f} = (f_i)_{i \in I}\) be a finite family of morphisms of varieties \(f_i: \mathbf{X}_i \to \mathbf{S}\) , and let \((\mathrm{J}_{\lambda})_{\lambda \in \Lambda}\) be a partition of I. Suppose that, for every \(\lambda\) in \(\Lambda\) , the family \(\mathbf{f}_{\lambda} = (f_i)_{i \in \mathrm{J}_{\lambda}}\) is transversal, and we denote by \(\mathrm{P}_{\lambda}\) the fiber product of this family; for every point \(x = (x_i)_{i \in \mathrm{J}_{\lambda}}\) of \(\mathrm{P}_{\lambda}\) , the element \(f_i(x_i)\) of S is independent of \(i \in \mathrm{J}_{\lambda}\) and will be denoted by \(u_{\lambda}(x)\) ;

then \(u_{\lambda}\) is a morphism from \(P_{\lambda}\) into S. For the family \(\mathbf{u} = (u_{\lambda})\) to be transversal, it is necessary and sufficient that the family \(\mathbf{f}\) be so. The canonical bijection of \(\prod_{\lambda \in \Lambda} \prod_{i \in J_{\lambda}} X_i\) onto \(\prod_{i \in I} X_i\) then gives by restriction an isomorphism from the fibered product \(\prod_{\lambda \in \Lambda} {}_{S}P_{\lambda}\) to the fibered product \(\prod_{i \in I} {}_{S}X_i\) .

5.11.5. Let S be a variety. One calls a variety over S a variety X equipped with a morphism \(\lambda: X \to S\) . Let \((X, \lambda)\) be a variety over S and \(f: S' \to S\) a morphism of varieties such that the pair \((f, \lambda)\) is transversal. We will then denote \(f^*(X)\) the variety \(S' \times_{S} X\) , endowed with the morphism \(f^*(\lambda): S' \times_{S} X\) defined by \(f^*(\lambda)(s', x) = s'\) . We say that \(f^*(X)\) is deduced from X by base change from S to \(S'\) along \(f\) . If \(\lambda\) is a submersion (resp. an immersion, a subimmersion, an étale morphism), the same holds for \(f^*(\lambda)\) .

5.11.6. Let \(f: X \to Y\) be a morphism of varieties and let W be a subvariety of Y; let \(i\) be the canonical injection of W into Y. We say that \(f\) is transverse to W at a point \(x \in f^{-1}(W)\) if the pair \((f, i)\) is transversal at the point \((x, f(x))\) of \(X \times W\) . For this, it is necessary and sufficient that the following conditions be satisfied:

\begin{enumerate}
\def\labelenumi{(\roman{enumi})}
\item
  the tangent space \(\mathbf{T}_{f(x)}(\mathbf{Y})\) is the sum of \(\mathbf{T}_{f(x)}(\mathbf{W})\) and the image of \(\mathbf{T}_{x}(f)\) ;
\item
  the inverse image \(\mathrm{T}_{x}(f)^{-1}(\mathrm{T}_{f(x)}(\mathrm{W}))\) of the tangent space to W at \(f(x)\) admits a topological complement.
\end{enumerate}

One says that \(f\) is transverse to \(W\) if it is transverse to \(W\) at every point of \(f^{-1}(W)\) .

For a morphism from X to Y to be a submersion, it suffices that it be transversal to every point of Y, and it is necessary that it be transversal to every submanifold of Y. For a finite family of morphisms \(f_{i}:X_{i}\to S\) to be transversal, it is necessary and sufficient that the morphism \(g\) from \(\prod_{i\in I}X_{i}\) to \(S^{I}\) defined by \(g((x_{i})_{i\in I})=(f_{i}(x_{i}))_{i\in I}\) be transverse to the diagonal of \(S^{I}\) .

5.11.7. Suppose that the morphism \(f:X\to Y\) is transverse to the subvariety W of Y. Then \(f^{-1}(W)\) is a submanifold of X, and for \(x\) in \(f^{-1}(W)\) the subspace of \(\mathrm{T}_{x}(\mathrm{X})\) tangent to \(f^{-1}(W)\) is the inverse image under \(\mathrm{T}_{x}(f)\) of the subspace \(\mathrm{T}_{f(x)}(\mathrm{W})\) . By passing to the quotient, the linear map \(\mathrm{T}_{x}(f)\) defines an isomorphism of topological vector spaces from the transverse space to \(f^{-1}(W)\) at \(x\) onto the transverse space to W at \(y=f(x)\) . If W has codimension \(d\) at \(y\) in Y, the submanifold \(f^{-1}(W)\) of X has codimension \(d\) at every point of \(f^{-1}(W)\) . Finally, the map \(x\mapsto(x,f(x))\) is an isomorphism of varieties from \(f^{-1}(W)\) to the fibered product \(\mathrm{X}\times_{\mathrm{Y}}\mathrm{W}\) .

5.11.8. Let \(Y_{1}\) and \(Y_{2}\) be two submanifolds of a variety X, and let \(\iota_{j}\) be the injection of \(Y_{j}\) into X. One says that \(Y_{1}\) and \(Y_{2}\) are transverse if the pair \((\iota_{1},\iota_{2})\) is transverse; this is equivalent to supposing that, for every point x of \(Y_{1}\cap Y_{2}\) , the subspaces \(T_{x}(Y_{1})\) and \(T_{x}(Y_{2})\) of \(T_{x}(X)\) have \(T_{x}(X)\) as their sum, and that their intersection admits a topological supplement in \(T_{x}(X)\) . Under these conditions, \(Y_{1}\cap Y_{2}\) is a submanifold of X and for every x in \(Y_{1}\cap Y_{2}\) , we have:

\[
\mathrm{T} _ {x} (\mathrm{Y} _ {1} \cap \mathrm{Y} _ {2}) = \mathrm{T} _ {x} (\mathrm{Y} _ {1}) \cap \mathrm{T} _ {x} (\mathrm{Y} _ {2});
\]

if moreover \(\mathbf{Y}_i\) has codimension \(d_i\) at \(x\) , then \(\mathbf{Y}_1 \cap \mathbf{Y}_2\) has codimension \(d_1 + d_2\) at \(x\) .

5.11.9. Let \(f\) and \(g\) be two morphisms from a variety X to a variety Y. If the morphism \((f,g):\mathrm{X}\to\mathrm{Y}\times\mathrm{Y}\) is transverse to the diagonal of \(\mathrm{Y}\times\mathrm{Y}\) , the subset N of X consisting of the points \(x\) such that \(f(x)=g(x)\) is a submanifold of X; it is called the kernel of the double arrow.

\[
f, g: \mathrm{X} \longrightarrow \mathrm{Y}.
\]

\subsection{Varieties of groups}

5.12.1. Let G be a group. A manifold structure on G is said to be compatible with the group structure of G if the map \((x, y) \mapsto xy\) of \(G \times G\) into G is a morphism. The map \(x \mapsto x^{-1}\) is then a morphism from G to itself. The set G, equipped with its group structure and its manifold structure, is called a group manifold (“of class \(C^{r}\)” if one wishes to be precise), or again a Lie group. If \(r = \omega\) we also say analytic group. If K = R (resp. C, \(Q_{p}\) ), one also says real (resp. complex, p-adic) Lie group. Every group variety is pure. A homomorphism of group varieties (or simply homomorphism) is any map from one group variety to another that is both a group homomorphism and a morphism of varieties.

If G is a group variety, the topological structure underlying the variety structure of G makes it a metrizable and complete topological group; it is locally compact if K is locally compact and G is of finite dimension.

5.12.2. Examples:

\begin{enumerate}
\def\labelenumi{(\roman{enumi})}
\item
  If V is a Banach space, the canonical manifold structure of V is compatible with its commutative group structure.
\item
  Let A be a complete normed K-algebra with a unit element, and let A* be the group of invertible elements of A. It is an open subspace of A, and the manifold structure induced on this open set by the canonical manifold structure of the K-vector space A is compatible with the group structure on A*. In particular, taking A to be the algebra \(\mathcal{L}(E)\) of continuous endomorphisms of a Banach space E; the group \(A^{*}\) consists of the automorphisms of the topological vector space E; we denote by \(\mathbf{GL}(E)\) the group variety thus defined. When \(A = M_{n}(K)\) , we obtain a group variety structure on \(\mathbf{GL}(n, K)\) .
\item
  If \(G_{1},\ldots,G_{n}\) are group varieties, the product group \(G=G_{1}\times\cdots\times G_{n}\) is a group variety, when endowed with the product variety structure of those of \(G_{i}\) .
\end{enumerate}

5.12.3. Let G be a group variety, let H be a topological group, and let \(f:H\to G\) be a continuous homomorphism satisfying condition (QR) of No. 5.8.1. The inverse image manifold structure of that of G by \(f\) then makes H a group variety. This applies in particular when H is a subgroup of G which is a subvariety (resp. quasi-subvariety); such a subgroup is called a group subvariety (resp. group quasi-subvariety) of G; it is necessarily closed in G.

If \(H_{i}(i=1,\ldots,n)\) is a group subvariety of \(G_{i}\) , then \(H_{1}\times\cdots\times H_{n}\) is a group subvariety of \(G_{1}\times\cdots\times G_{n}\) .

5.12.4. Let G be a group variety and let H be a group subvariety of G. The equivalence relation \(x^{-1}y \in H\) is regular, which allows one to endow the space G/H of left classes xH with a variety structure called the quotient of that of G. The canonical map \((g, x) \mapsto g.x\) of \(G \times (G/H)\) into G/H is a morphism. Analogous results hold for the homogeneous space \(H\backslash G\) of right classes Hx. If H is normal, the variety structure of G/H is compatible with its group structure.

5.12.5. Let G be a group variety and let X be a variety. A left action of the group G on the variety X is a morphism \((s, x) \mapsto sx\) of \(G \times X\) into X such that

\[
\begin{array}{l l} s (t x) = (s t) x & \text { if } \quad s, t \in G, x \in X \\ e x = x & \text { if } \quad x \in X (e \text { being   the   identity   element   of } G). \end{array}
\]

One also says that the group variety G operates on the left on X; right action laws are defined analogously.

Let \(x \in X\) and let \(G_x\) be its stabilizer in \(G\) . Suppose that the map \(g \mapsto g \cdot x\) is a subimmersion (a hypothesis automatically satisfied if the characteristic of \(K\) is 0 and \(X\) is of finite dimension). Then \(G_x\) is a group subvariety of \(G\) and the map from \(G/G_x\) to \(X\) obtained by passing to the quotient from \(g \mapsto g \cdot x\) is an immersion. If moreover the orbit \(G \cdot x\) of \(x\) is locally closed and if the topology of \(G\) admits a countable base, \(G \cdot x\) is a subvariety of \(X\) and the map \(G/G_x \to G \cdot x\) is an isomorphism of varieties.

\subsection{Weakening of structure}

Throughout this section, the letters r, s, r', s' denote either integers \(\geqslant1\) , or one of the symbols \(\infty\) and \(\omega\) . One assumes that K = R.

5.13.1. Let \(r \leqslant s\) , and X be a variety of class \(C^s\) . There exists on the topological space X a manifold structure of class \(C^r\) and only one such that every chart of X for the given structure is a chart of X for this new structure. Let \(X_r\) be the manifold of class \(C^r\) thus defined. It is said to be deduced from X by weakening the manifold structure of X, or again that its manifold structure is underlying that of X. The notion of weakening is transitive: if \(r' \leqslant r\) , one has \(X_{r'} = (X_r)_{r'}\) ; it commutes with products: if Y is of class \(C^s\) , one has \((X \times Y)_r = X_r \times Y_r\) ; one has an analogous result for \(X \times_s Y\) under the hypotheses of no. 5.11.2.

Let \(a\in X\) and let \(c\) be a chart of \(X\) at \(a\) ; \(c\) is also a chart of \(X_r\) at \(a\) . One deduces from this (5.5.1) the isomorphisms

\[
\theta_ {c}: \mathrm{E} \rightarrow \mathrm{T} _ {a} (\mathrm{X}), \quad \theta_ {c} ^ {\prime}: \mathrm{E} \rightarrow \mathrm{T} _ {a} (\mathrm{X} _ {r}),
\]

whence an isomorphism \(\theta_{c}^{\prime}\circ\theta_{c}^{-1}:T_{a}(X)\to T_{a}(X_{r})\) . This isomorphism is independent of the choice of \(c\) ; it allows one to identify the tangent spaces to X and to \(X_{r}\) at \(a\) .

5.13.2. Let X (resp. \(\mathbf{X}'\) ) be a manifold of class \(\mathbf{C}^s\) (resp. \(\mathbf{C}^{s'}\) ) and let \(r\) satisfy \(r \leqslant \inf(s, s')\) . A map \(f: \mathbf{X} \to \mathbf{X}'\) is said to be of class \(\mathbf{C}^r\) if it is a morphism of \(\mathbf{X}_r\) to \(\mathbf{X}_r'\) ; such a map is of class \(\mathbf{C}^{r'}\) for all \(r' \leqslant r\) . Moreover, the tangent linear map to \(f: \mathbf{X}_r \to \mathbf{X}_r'\) at a point \(a \in X\) coincides with the tangent linear map to \(f\) considered as a morphism of \(\mathbf{X}_{r'}\) to \(\mathbf{X}_{r'}'\) .

One will most often denote by the same symbol the manifolds X and \(X_{r}\) ; thus, if X is of class \(C^{s}\) , the expression “a submanifold of X of class \(C^{r}\)” ( \(r \leqslant s\) ) means “a submanifold of \(X_{r}\)”.

\subsection{Restriction of the base field}

In this number, one is given two complete non-discrete valued commutative fields K and L, as well as an isomorphism \(\sigma\) of the valued field K onto a subfield of L. If E is a Banach space over L, one denotes by \(\sigma_{*}(\mathbf{E})\) the vector space over K obtained by restriction of scalars (cf.~Alg., chap.~II, 3rd ed., § 8); the given topology on E is compatible with the structure of K-vector space, and \(\sigma_{*}(\mathbf{E})\) is a Banach space over K.

5.14.1. Let X be an analytic manifold over L and let \(c = (U, \varphi, E)\) be a chart of X. The map \(\varphi\) is a bijection of U onto an open subset of \(\sigma_{*}(E)\) and the triple \(c_{\sigma} = (U, \varphi, \sigma_{*}(E))\) is a chart of X. There exists on X a structure of analytic manifold over K and only one for which \(c_{\sigma}\) is a chart for every chart c of the L-analytic manifold X. One denotes by \(X_{\sigma}\) the analytic manifold over K thus obtained and one says that \(X_{\sigma}\) is obtained from X by restriction of scalars (from L to K by means of \(\sigma\) ). The topological space underlying \(X_{\sigma}\) is the same as that underlying X.

5.14.2. Examples:

\begin{enumerate}
\def\labelenumi{(\alph{enumi})}
\item
  One takes K = R, L = C, σ being the canonical injection of R into C. Thus, every complex analytic manifold is canonically equipped with a structure of real analytic manifold; this real analytic structure itself defines differential structures of class \(C^{r}\) for every r.
\item
  One takes K = C, L = C, σ being the conjugation \(x \mapsto \bar{x}\) . One thus associates to every complex analytic manifold X a complex analytic manifold \(\bar{X}\) , called the conjugate of X. If f is a complex-valued function, defined on an open subset U of X, f is analytic for the structure of \(\bar{X}\) if and only if the conjugate function \(\bar{f}\) is analytic for the structure of X. The manifolds X and \(\bar{X}\) define by restriction of scalars the same real analytic manifold.
\end{enumerate}

5.14.3. Let X be an analytic manifold over K, Y an analytic manifold over L. A map \(f:X\to Y\) is said to be \(\sigma\) -analytic if it is a morphism of X into \(Y_{\sigma}\) . If \(K\subset L\) , \(\sigma\) is the injection of K into L, and X is an analytic manifold over L, one calls K-analytic map a map \(\sigma\) -analytic from \(X_{\sigma}\) into Y.

5.14.4. Let V be a Banach space over L. One has \(V_{\sigma} = \sigma_{*}(V)\) ; the canonical analytic variety structure over K of \(\sigma_{*}(V)\) is deduced by restriction of scalars from the canonical analytic variety structure over L of V.

5.14.5. Let X be an analytic variety over L, and let \(a \in X\) . Let c be a chart of X at a; then \(c_{\sigma}\) is a chart of \(X_{\sigma}\) at a, and one deduces from it isomorphisms

\[
\theta_ {c} \colon \mathrm{E} \rightarrow \mathrm{T} _ {a} (\mathrm{X}), \quad \theta_ {c _ {\sigma}} \colon \sigma_ {*} (\mathrm{E}) \rightarrow \mathrm{T} _ {a} (\mathrm{X} _ {\sigma})
\]

whence an isomorphism \(\theta_{c_{\sigma}} \circ \sigma_{*}(\theta_{c})^{-1}\) of \(\sigma_{*}(\mathrm{T}_{a}(\mathrm{X}))\) onto \(\mathrm{T}_{a}(\mathrm{X}_{\sigma})\) ; this isomorphism does not depend on the choice of \(c\) ; it allows one to identify \(\mathrm{T}_{a}(\mathrm{X}_{\sigma})\) with \(\sigma_{*}(\mathrm{T}_{a}(\mathrm{X}))\) and even with \(\mathrm{T}_{a}(\mathrm{X})\) by abuse of notation.

If \(f\) is an analytic map on L from X into a manifold Y, the tangent linear map to \(f\) at \(a\) ( \(f\) being regarded as a morphism of \(\mathrm{X}_{\sigma}\) to \(\mathrm{Y}_{\sigma}\) ) is equal to \(\sigma_{*}(\mathrm{T}_{a}(f))\) .

5.14.6. Let X and Y be two analytic varieties over L, and let \(f: \mathbf{X}_{\sigma} \to \mathbf{Y}\) be a continuous \(\sigma\) -analytic map. Suppose that the characteristic of K is 0. Then, for \(f\) to be analytic over L, it is necessary and sufficient that, for every \(a \in \mathbf{X}\) , the map \(\mathrm{T}_{a}(f)\) be L-linear.

When K = R, L = C (case (a) of no. 5.14.2), one has a more precise result: if X and Y are finite-dimensional and if \(f: X \to Y\) is a map of class \(C^{1}\) whose tangent map at every point of X is C-linear, then \(f\) is complex analytic.

5.14.7. Let X be a complex analytic manifold and g a map from X into itself satisfying the conditions:

\begin{enumerate}
\def\labelenumi{(\roman{enumi})}
\item
  One has \(g \circ g = \mathrm{Id}_{\mathbf{X}}\) .
\item
  The map \(g\) is an isomorphism of the analytic variety \(X\) on the conjugate variety \(\overline{X}\) (5.14.2).
\end{enumerate}

The set \(X_{0}\) of points x of X such that \(g(x) = x\) is a closed analytic submanifold of the real analytic manifold underlying X. For \(x \in X_{0}\) , one has \(\mathbf{T}_{x}(\mathbf{X}) = \mathbf{T}_{x}(\mathbf{X}_{0}) \oplus i\mathbf{T}_{x}(\mathbf{X}_{0})\) .

Let U be a connected open subset of X, and let f and g be two complex analytic maps from U into a complex separated locally convex space or into a complex separated analytic manifold. If f and g coincide on a nonempty subset of \(U \cap X_{0}\) , open in \(X_{0}\) , then f = g.

Suppose \(X_{0}\) is paracompact. If f is a real analytic map from \(X_{0}\) into a separated locally convex space or into a separated complex analytic manifold, there exists an open neighborhood U of \(X_{0}\) in X and a complex analytic map from U into the space of values of f, extending f. Any two such extensions coincide on a neighborhood of \(X_{0}\) in X.

Suppose X is of finite dimension. For every point \(a \in X_{0}\) , there exists a system of (complex) coordinates \(\zeta^{1}, \ldots, \zeta^{n}\) in an open neighborhood U of \(a\) , such that \(\zeta^{i} \circ g = \bar{\zeta}^{i}\) for \(1 \leqslant i \leqslant n\) ; the restriction \(\xi^{i}\) of \(\zeta^{i}\) to \(U \cap X_{0}\) then takes real values and \((\xi^{1}, \ldots, \xi^{n})\) is a system of coordinates at \(a\) of the real analytic manifold \(X_{0}\) .

5.14.8. For every paracompact real analytic manifold Y, there exists a pair \((X, g)\) consisting of a complex analytic manifold X and a map g from X into itself satisfying conditions (i) and (ii) of 5.14.7, and an isomorphism f from Y onto \(X_{0}\) . We say that X (endowed with f and g) is a complexification of Y.

\section{\texorpdfstring{Fibrations \(^{1}\)}{Fibrations ¹}}

\subsection{Fibrations}

6.1.1. A fibration of class \(C^{r}\) , or simply fibration, is a triple \((X, B, \pi)\) where B and X are varieties of class \(C^{r}\) and \(\pi\) is a morphism from X to B, having the following property:

\begin{enumerate}
\def\labelenumi{(\Alph{enumi})}
\setcounter{enumi}{5}
\tightlist
\item
  For every \(x \in \mathbf{B}\) , there exists an open neighborhood \(\mathbf{U}\) of \(x\) , a manifold \(\mathbf{F}\) , and an isomorphism \(\varphi\) of \(\pi^{-1}(\mathbf{U})\) onto \(\mathbf{U} \times \mathbf{F}\) such that \(\pi(\varphi^{-1}(x, y)) = x\) for all \(x \in \mathbf{U}\) and every \(y \in \mathbf{F}\) .
\end{enumerate}

If \(\lambda = (X, B, \pi)\) is a fibration, one calls \(X\) the space of \(\lambda\) , \(B\) the base of \(\lambda\) , and \(\pi\) the projection of \(\lambda\) . The map \(\pi\) is a submersion; in particular \(\pi(X)\) is open in \(B\) and if \(R\) denotes the equivalence relation defined by \(\pi\) on \(X\) , the canonical map from \(X/R\) to \(\pi(X)\) is an isomorphism. For every \(x \in B\) the inverse image \(\pi^{-1}(x)\) is a closed submanifold of \(X\) , called the fiber of \(x\) , and denoted \(X_x\) .

6.1.2. Examples:

\begin{enumerate}
\def\labelenumi{(\alph{enumi})}
\item
  If B and F are two manifolds, the triple \((\mathbf{B} \times \mathbf{F}, \mathbf{B}, \mathbf{pr}_{1})\) is a fibration whose fibers are canonically isomorphic to F.
\item
  If \(\lambda = (\mathbf{X},\mathbf{B},\pi)\) and \(\lambda^{\prime} = (\mathbf{X}^{\prime},\mathbf{B}^{\prime},\pi^{\prime})\) are fibrations,
\end{enumerate}

\[
\lambda \times \lambda^ {\prime} = (\mathrm{X} \times \mathrm{X} ^ {\prime}, \mathrm{B} \times \mathrm{B} ^ {\prime}, \pi \times \pi^ {\prime})
\]

is a fibration; it is called the product of \(\lambda\) and of \(\lambda'\) .

\begin{enumerate}
\def\labelenumi{(\alph{enumi})}
\setcounter{enumi}{2}
\tightlist
\item
  If \(\lambda = (X, B, \pi)\) and \(\lambda' = (X', B, \pi')\) are fibrations with the same base, \(\lambda \times_{B} \lambda' = (X \times_{B} X', B, \pi \times_{B} \pi')\) is a fibration; it is called the product of \(\lambda\) and of \(\lambda'\) above \(B\) , or also the fibered product of \(\lambda\) and of \(\lambda'\) .
\end{enumerate}

6.1.3. Let \(\lambda = (X, B, \pi)\) and \(\lambda' = (X', B', \pi')\) be two fibrations. One calls a morphism of \(\lambda\) to \(\lambda'\) any pair \((f, g)\) , where \(f\) is a morphism of \(B\) to \(B'\) , and \(g\) a morphism of \(X\) to \(X'\) such that \(\pi' \circ g = f \circ \pi\) . When \(B = B'\) and \(f = \mathrm{Id}_B\) , one says that \(g\) is a \(B\) -morphism of \(\lambda\) to \(\lambda'\) ; if \(g\) is an isomorphism of \(X\) onto \(X'\) , \(g^{-1}\) is a \(B\) -morphism of \(\lambda'\) to \(\lambda\) , and one says that \(g\) is a \(B\) -isomorphism from \(\lambda\) onto \(\lambda'\) ; for this to be the case, it is necessary and sufficient that, for every \(x \in B\) , the map \(g_x: X_x \to X'_x\) induced by \(g\) be an isomorphism.

6.1.4. The fibration \((\mathbf{B} \times \mathbf{F}, \mathbf{B}, \mathbf{pr}_{1})\) is called the trivial fibration with base B and fiber F. An isomorphism of a fibration \(\lambda\) onto a trivial fibration is called a trivialization of \(\lambda\) .

6.1.5. Let \(\lambda = (X, B, \pi)\) be a fibration, and let \(f: B' \to B\) be a morphism. Let \(\pi'\) be the canonical morphism of \(B' \times_B X\) to \(B'\) . The triple \((B' \times_B X, B', \pi')\) is a fibration, called the inverse image of \(\lambda\) by \(f\) or the fibration deduced from \(\lambda\) by base change of \(B\) to \(B'\) along \(f\) , and denoted \(B' \times_B \lambda\) , or again \(f^*\lambda\) . If \(f'\) denotes the canonical map of \(B' \times_B X\) to \(X\) , the pair \((f, f')\) is a morphism of \(B' \times_B \lambda\) to \(\lambda\) ; it enjoys the following universal property: if \((f, g)\) is a morphism of a fibration \(\lambda'\) with base \(B'\) into the fibration \(\lambda\) , there exists a unique \(B'\) -morphism \(\varphi: \lambda' \to B' \times_B \lambda\) such that \((f, g) = (f, f') \circ \varphi\) .

When B' is a submanifold of B, and f is the canonical injection of B' into B, \(B' \times_{B} X\) is identified with the submanifold \(\pi^{-1}(B')\) of X, and \(\pi'\) is identified with the restriction of \(\pi\) to \(\pi^{-1}(B')\) ; the inverse image of \(\lambda\) by f is then called the fibration induced by \(\lambda\) on B'.

6.1.6. If \(\lambda = (X, B, \pi)\) is a fibration, one calls a morphic section (or simply section) of \(\lambda\) every morphism \(s: B \to X\) such that \(\pi \circ s = \mathrm{Id}_B\) .

\subsection{Principal fibrations}

6.2.1. Let B be a manifold and G a group manifold. A principal fibration with base B and structure group G is called a quadruple \(\lambda = (P, G, B, \pi)\) where P is a variety on which G acts on the right by \((x, g) \mapsto x \cdot g\) (cf. no. 5.12.5), and where \(\pi\) is a morphism from P to B, these data being subject to the following axiom:

\begin{enumerate}
\def\labelenumi{(\Alph{enumi})}
\setcounter{enumi}{15}
\tightlist
\item
  For every \(b \in \mathbf{B}\) there exists an open neighborhood \(\mathbf{U}\) of \(b\) and an isomorphism \(f: \mathbf{U} \times \mathbf{G} \to \pi^{-1}(\mathbf{U})\) such that we have
\end{enumerate}

\[
\pi (f (u, g)) = u \text {   and   } f (u, g g ^ {\prime}) = f (u, g) \cdot g ^ {\prime} \text {   if   } u \in \mathbf {U} \text {   and   } g, g ^ {\prime} \in \mathbf {G}.
\]

6.2.2. Let \(\lambda = (\mathbf{P}, \mathbf{G}, \mathbf{B}, \pi)\) be a principal fibration. The triple \((\mathbf{P}, \mathbf{B}, \pi)\) is a fibration; we have \(\pi(\mathbf{P}) = \mathbf{B}\) . The equivalence relation \(\mathbf{R}\) defined by \(\pi\) on \(\mathbf{P}\) coincides with the one defined by \(\mathbf{G}\) ; its graph is none other than the product \(\mathbf{P} \times_{\mathbf{B}} \mathbf{P}\) (cf.~no. 5.11.2); it is a submanifold of \(\mathbf{P} \times \mathbf{P}\) . The map \((x, g) \mapsto (x, x \cdot g)\) is an isomorphism of \(\mathbf{P} \times \mathbf{G}\) onto \(\mathbf{P} \times_{\mathbf{B}} \mathbf{P}\) ; the map which assigns to every \((x, y) \in \mathbf{P} \times_{\mathbf{B}} \mathbf{P}\) the unique element \(g \in \mathbf{G}\) such that \(y = x \cdot g\) is a morphism of \(\mathbf{P} \times_{\mathbf{B}} \mathbf{P}\) to \(\mathbf{G}\) .

The group G acts properly and freely on P (cf. Gen. Top., chap. III, 3rd ed., § 4). If \(x \in \mathbf{P}\) and if \(b = \pi(x)\) , the map \(g \mapsto x \cdot g\) is an isomorphism from the variety G onto the fiber over b.

6.2.3. Conversely, let G be a group variety, and let P be a variety on which G acts on the right so as to satisfy the following two conditions:

\begin{enumerate}
\def\labelenumi{(\alph{enumi})}
\item
  G acts properly and freely on P.
\item
  For every \(x \in P\) , the map \(g \mapsto x \cdot g\) is an immersion of G into P.
\end{enumerate}

Then the equivalence relation defined by G in P is regular; if P/G denotes the quotient variety, and \(\pi\) the canonical projection from P to P/G, the quadruple (P, G, P/G, \(\pi\) ) is a principal fibration.

When K has characteristic 0 and P is finite-dimensional, condition (b) above is a consequence of condition (a).

6.2.4. The conditions of the preceding number are satisfied in the following two cases:

\begin{enumerate}
\def\labelenumi{(\alph{enumi})}
\item
  P is a group variety and G a subgroup variety acting on P by right translations; the base of the principal fibration thus obtained is the homogeneous space P/G.
\item
  G is a discrete group acting properly and freely on P. The projection \(\pi: \mathrm{P} \to \mathrm{P} / \mathrm{G}\) is then an étale morphism (no. 5.7.6).
\end{enumerate}

6.2.5. Examples:

\begin{enumerate}
\def\labelenumi{(\alph{enumi})}
\item
  Let B be a variety and let G be a group variety. Let G act on \(B \times G\) by \((b, g) \cdot g' = (b, gg')\) . The quadruple \((B \times G, G, B, pr_{1})\) is a principal fibration.
\item
  Let \(\lambda = (\mathbf{P}, \mathbf{G}, \mathbf{B}, \pi)\) and \(\lambda' = (\mathbf{P}', \mathbf{G}', \mathbf{B}', \pi')\) be two principal fibrations. Let us make \(\mathbf{G} \times \mathbf{G}'\) act on \(\mathbf{P} \times \mathbf{P}'\) by the formula:
\end{enumerate}

\[
(x, x ^ {\prime}). (g, g ^ {\prime}) = (x. g, x ^ {\prime}. g ^ {\prime}), \quad x \in \mathbf {P}, x ^ {\prime} \in \mathbf {P} ^ {\prime}, g \in \mathbf {G}, g ^ {\prime} \in \mathbf {G} ^ {\prime}.
\]

The quadruple \(\lambda \times \lambda' = (\mathbf{P} \times \mathbf{P}', \mathbf{G} \times \mathbf{G}', \mathbf{B} \times \mathbf{B}', \pi \times \pi')\) is a principal fibration; it is called the product of \(\lambda\) and of \(\lambda'\) .

\begin{enumerate}
\def\labelenumi{(\alph{enumi})}
\setcounter{enumi}{2}
\tightlist
\item
  Let \(\lambda = (P, G, B, \pi)\) and \(\lambda' = (P', G', B, \pi')\) two principal fibrations with the same base. The submanifold \(P \times_{B} P'\) of \(P \times P'\) is stable under the operations of \(G \times G'\) , and the quadruple
\end{enumerate}

\[
\lambda \times_ {\mathbf {B}} \lambda^ {\prime} = (\mathbf {P} \times_ {\mathbf {B}} \mathbf {P} ^ {\prime}, \mathbf {G} \times \mathbf {G} ^ {\prime}, \mathbf {B}, \pi \times_ {\mathbf {B}} \pi^ {\prime})
\]

is a principal fibration; it is called the product of \(\lambda\) and of \(\lambda'\) above B, or equivalently the fiber product of \(\lambda\) and of \(\lambda'\) .

\subsection{Morphisms of principal fibrations}

6.3.1. Let \(\lambda = (\mathbf{P}, \mathbf{G}, \mathbf{B}, \pi)\) and \(\lambda' = (\mathbf{P}', \mathbf{G}', \mathbf{B}', \pi')\) be two principal fibrations. A morphism of \(\lambda\) to \(\lambda'\) is a triple \((f, \varphi, h)\) , where \(f: \mathbf{P} \to \mathbf{P}'\) and \(h: \mathbf{B} \to \mathbf{B}'\) are morphisms, \(\varphi: \mathbf{G} \to \mathbf{G}'\) is a homomorphism of group varieties, and where \(\pi' \circ f = h \circ \pi\) and \(f(x \cdot g) = f(x) \cdot \varphi(g)\)

for \(x \in \mathbf{P}, g \in \mathbf{G}\) . Note that \(f\) determines \(h\) ; one will often say that \((f, \varphi)\) , or even simply \(f\) is a morphism.

When \(B = B'\) and \(h = \mathrm{Id}_{B}\) (resp. when \(G = G'\) and \(\varphi = \mathrm{Id}_{G}\) ), a morphism is called a B-morphism compatible with \(\varphi\) (resp. a G-morphism compatible with h). A morphism that is both a B-morphism and a G-morphism is called a G-B-morphism (here again, one often says "morphism" when there can be no confusion). Every G-B-morphism \(f : P \to P'\) is an isomorphism of the variety P onto the variety \(P'\) ; the inverse isomorphism \(f^{-1}\) is a G-B-morphism: f is a G-B-isomorphism of P onto \(P'\) .

Two principal fiber bundles P and P' with the same base B and the same structural group G are said to be G-B-isomorphic (or simply isomorphic) if there exists a G-B-isomorphism of P onto P'.

6.3.2. The principal fibration \((\mathbf{B} \times \mathbf{G}, \mathbf{G}, \mathbf{B}, \mathbf{pr}_{1})\) is called the trivial principal fibration with base B and structural group G. An isomorphism of a principal fibration \(\lambda = (\mathbf{P}, \mathbf{G}, \mathbf{B}, \pi)\) onto the trivial principal fibration with base B and structural group G is called a trivialization of \(\lambda\) . Every section s of \(\lambda\) defines a trivialization \(f_{s}\) of \(\lambda\) by the formula:

\[
f _ {s} ^ {- 1} (b, g) = s (b). g \quad \text {   for   } b \in \mathbf {B} \text {   and   } g \in \mathbf {G}.
\]

Thus we obtain a bijection from the set of sections of \(\lambda\) onto the set of trivializations of \(\lambda\) . Moreover, if \(s_0\) is a section of \(\lambda\) , every section \(s\) of \(\lambda\) can be written uniquely in the form \(s(b) = s_0(b) \cdot r(b)\) , where \(r: B \to G\) is a morphism.

6.3.3. Let \(\lambda = (\mathbf{P}, \mathbf{G}, \mathbf{B}, \pi)\) be a principal fibration, and let \(h: \mathbf{B}' \to \mathbf{B}\) be a morphism. Let \(\pi'\) (resp. \(h'\) ) be the canonical morphism from \(\mathbf{B}' \times_{\mathbf{B}} \mathbf{P}\) to \(\mathbf{B}'\) (resp. into \(\mathbf{P}\) ). Let us make \(\mathbf{G}\) act on \(\mathbf{B}' \times_{\mathbf{B}} \mathbf{P}\) by the formula

\[
(b ^ {\prime}, x). g = (b ^ {\prime}, x. g), \quad (b ^ {\prime}, x) \in \mathbf {B} ^ {\prime} \times_ {\mathbf {B}} \mathbf {P}, \quad g \in \mathbf {G}.
\]

The quadruple \((\mathbf{B}' \times_{\mathbf{B}} \mathbf{P}, \mathbf{G}, \mathbf{B}', \pi')\) is a principal fibration, called the inverse image of \(\lambda\) by \(h\) , and denoted \(\mathbf{B}' \times_{\mathbf{B}} \lambda\) or also \(h^*\lambda\) . The map \(h': \mathbf{B}' \times_{\mathbf{B}} \mathbf{P} \to \mathbf{P}\) is a G-morphism compatible with \(h\) ; it enjoys the following universal property: if \(f\) is a G-morphism compatible with \(h\) of a principal fibration \(\lambda'\) with base \(\mathbf{B}'\) into the fibration \(\lambda\) , there exists a unique G-B'-isomorphism \(k: \lambda' \to \mathbf{B}' \times_{\mathbf{B}} \lambda\) such that \(f = h' \circ k\) .

When B' is a subvariety of B, and h is the canonical injection of B' into B, \(B' \times_{B} P\) is identified with the submanifold \(\pi^{-1}(B')\) of P; it is called the principal fiber bundle induced by P over B', and it is denoted \(\pi^{-1}(B')\) , or equivalently P\textbar B'. Every \(x \in B\) has an open neighborhood U such that P\textbar U is trivial.

\subsection{Construction of principal fibrations by means of cocycles}

Let B be a manifold, G a group manifold, and let \(\mathcal{U} = (\mathrm{U}_{i})_{i\in I}\) an open cover of B.

6.4.1. One calls a cocycle of class \(C^{r}\) on B with values in G, subordinate to U, a family \((g_{i,j})_{(i,j)\in\mathbf{I}\times\mathbf{I}}\) possessing the following two properties: (1) for every pair \((i,j)\in\mathbf{I}\times\mathbf{I}\) , \(g_{i,j}\) is a map of class \(C^{r}\) of the open subset \(U_{i}\cap U_{j}\) of B into G;

\begin{enumerate}
\def\labelenumi{(\arabic{enumi})}
\setcounter{enumi}{1}
\tightlist
\item
  for every triple \((i,j,k)\in \mathbf{I}^3\) , one has
\end{enumerate}

\[
g _ {i, k} (x) = g _ {i, j} (x). g _ {j, k} (x) \quad \text { for   every } \quad x \in \mathrm{U} _ {i} \cap \mathrm{U} _ {j} \cap \mathrm{U} _ {k}.
\]

Two such cocycles \((g_{i,j})\) and \((g'_{i,j})\) are said to be cohomologous if there exists a family \((h_{i})_{i\in I}\) where, for every \(i\in I\) , \(h_{i}\) is a map of class \(C^{r}\) of \(U_{i}\) into G, such that:

\[
g _ {i, j} ^ {\prime} (x) = h _ {i} (x) ^ {- 1}. g _ {i, j} (x). h _ {j} (x) \quad \text { for   every } \quad x \in \mathrm{U} _ {i} \cap \mathrm{U} _ {j}.\tag{3}
\]

6.4.2. Let \(\lambda = (\mathbf{P},\mathbf{G},\mathbf{B},\pi)\) be a principal fibration. For every \(i\in \mathbf{I}\) let \(s_i\) be a section of \(\lambda\) over \(\mathbf{U}_i\) (6.3.3). For every pair \((i,j)\in \mathrm{I}^2\) there then exists one and only one morphism \(g_{i,j}\) of \(\mathbf{U}_i\cap \mathbf{U}_j\) to \(\mathbf{G}\) such that:

\[
s _ {j} (b) = s _ {i} (b). g _ {i, j} (b) \quad \text { for   every } \quad b \in \mathrm{U} _ {i} \cap \mathrm{U} _ {j}.\tag{4}
\]

The family of \(g_{i,j}\) is a cocycle on B with values in G, subordinate to the open covering \(\mathcal{U}\) . This cocycle is said to be associated with the object \((\lambda, \mathcal{U}, (s_i)_{i \in I})\) and the maps \(g_{i,j}\) are called the transition functions of this object.

For \(i \in I\) let \(x \mapsto (\pi(x), f_i(x))\) be the trivialization defined by the section \(s_i\) of \(\lambda |U_i\) (6.3.2). For \(x \in \pi^{-1}(U_i \cap U_j)\) , we have:

\[
f _ {i} (x) = g _ {i, j} (\pi (x)). f _ {j} (x).\tag{5}
\]

6.4.3. Conversely, let \(g = (g_{i,j})\) be a cocycle on B with values in G, subordinate to the covering \(\mathcal{U}\) . There then exists a principal fibration \(\lambda = (\mathrm{P},\mathrm{G},\mathrm{B},\pi)\) and a family of sections \((s_i)_{i\in I}\) of \(\lambda\) above the \(U_i\) such that relation (4) is satisfied. Then (5) also holds. If moreover \((\lambda',(s_i'))\) satisfies the same conditions, there exists a unique G-B-isomorphism \(f\) of \(\lambda\) onto \(\lambda'\) such that \(s_i' = f\circ s_i\) for all \(i\in I\) . We express this result by saying that \((\lambda,(s_i))\) is determined up to unique isomorphism by the cocycle \(g\) .

6.4.4. Let \(\lambda = (\mathbf{P}, \mathbf{G}, \mathbf{B}, \pi)\) and \(\lambda' = (\mathbf{P}', \mathbf{G}, \mathbf{B}, \pi')\) be two principal fibrations. Let \((s_i)\) (resp. \((s_i')\) ) be a family of sections of \(\lambda\) (resp. \(\lambda'\) ) over the \(U_i\) and let \(g\) (resp. \(g'\) ) be the cocycle associated with \((\lambda, \mathcal{U}, (s_i))\) (resp. \((\lambda', \mathcal{U}, (s_i'))\) ). For \(\lambda\) and \(\lambda'\) to be G-B-isomorphic, it is necessary and sufficient that

the cocycles g and g′ be cohomologous. More precisely, for every G-B-isomorphism f from λ onto λ′, there exists a family \((h_{i})_{i\in I}\) of morphisms of \(U_{i}\) into G and exactly one such that relation (3) is satisfied and that we have \(f\circ s_{i}^{\prime}(x)=s_{i}(x) \cdot h_{i}(x)\) for all \(i\in I\) and all \(x\in U_{i}\) ; one thus obtains a bijection from the set of G-B-isomorphisms of λ onto λ' onto the set of families \((h_{i})_{i\in I}\) satisfying (3).

6.4.5. We return to the notations of 6.4.2. and let \(\mathcal{V} = (\mathrm{V}_{\alpha})_{\alpha \in \mathrm{A}}\) be an open covering finer than the open covering \(\mathcal{U}\) . Let \(\tau : \mathrm{A} \to \mathrm{I}\) be a map such that \(\mathrm{V}_{\alpha} \subset \mathrm{U}_{\tau(\alpha)}\) for all \(\alpha \in \mathrm{A}\) . Let \(s_{\alpha}'\) be the restriction to \(\mathrm{V}_{\alpha}\) of the section \(s_{\tau(\alpha)}\) and let \(g' = (g_{\alpha,\beta}')\) be the cocycle subordinate to the open covering \(\mathcal{V}\) associated with \((\lambda, \mathcal{V}, (s_{\alpha}'))\) . The transition function \(g_{\alpha,\beta}'\) is then the restriction to \(\mathrm{V}_{\alpha} \cap \mathrm{V}_{\beta}\) of the transition function \(g_{\tau(\alpha),\tau(\beta)}\) .

\subsection{Fiber bundles associated with a principal fibration}

6.5.1. Let \(\lambda = (\mathbf{P}, \mathbf{G}, \mathbf{B}, \pi)\) be a principal fibration. Let F be a variety on which the group G acts on the left; we denote \((g, y) \mapsto g \cdot y\) the operation law of G on F. The group G acts on the right on \(\mathbf{P} \times \mathbf{F}\) by the formula \((x, f) \cdot g = (x \cdot g, g^{-1} \cdot f)\) ; the equivalence relation defined by G in \(\mathbf{P} \times \mathbf{F}\) is regular; the quotient \(\mathbf{P} \times^{\mathbf{G}}\mathbf{F} = (\mathbf{P} \times \mathbf{F}) / \mathbf{G}\) is endowed with a manifold structure.

Let E be a manifold. We say that E is equipped with a structure of fiber bundle associated with λ with typical fiber F when we are given a morphism ρ : P × F → E satisfying the following property:

(As) We have \(\rho(x \cdot g,g^{-1} \cdot f)=\rho(x,f)\) for \(x\in\mathbf{P}\) , \(f\in\mathbf{F}\) , \(g\in\mathbf{G}\) , and the map \(\bar{\rho}:\mathbf{P}\times^{\mathbf{G}}\mathbf{F}\to\mathbf{E}\) deduced from \(\rho\) by passing to the quotient is an isomorphism of varieties.

It amounts to the same thing to say that \((\mathbf{P} \times \mathbf{F}, \mathbf{G}, \mathbf{E}, \rho)\) is a principal fibration. The map \(\rho\) (or sometimes the map \(\bar{\rho}\) ) is called the frame map of E, and is denoted \((x, f) \mapsto x \cdot f\) ; we have

\[
(x. g). f = x. (g. f), \quad \text {   for   } x \in \mathbf {P}, g \in \mathbf {G} \text {   and   } f \in \mathbf {F}.
\]

The data of \(\lambda\) and F determines E up to a unique isomorphism; in particular, one may take for E the variety P × \(^{\text{G}}\) F itself; it is said to be the associated fiber bundle to \(\lambda\) of fiber type F, and it is denoted \(\lambda(\text{F})\) .

6.5.2. Let E be a fiber bundle associated with \(\lambda\) and with typical fiber F. There exists a unique morphism \(\pi_{\mathbf{E}}\) from E into B such that \(\pi_{\mathbf{E}}(x,f)=\pi(x)\) if \(x\in\mathbf{P},f\in\mathbf{F}\) ; the triple (E, B, \(\pi_{\mathbf{E}}\) ) is a fibration; if B and F are separated, then E is separated.

Let \(b \in \mathbf{B}\) , and let \(\mathbf{F}_b = \pi_{\mathbf{E}}^{-1}(b)\) ; it is a closed submanifold of E. If \(x \in \mathbf{P}\) is such that \(\pi(x) = b\) , let \(\theta_x: \mathbf{F} \to \mathbf{F}_b\) be the map defined by \(\theta_x(f) = x.f\) ; it is an isomorphism of manifolds. Moreover, for every \(g \in \mathbf{G}\) , one has \(\theta_{x.g} = \theta_x \circ \rho_g\) , where \(\rho_g\) denotes the automorphism \(f \mapsto g.f\) of F. Suppose

F is equipped with a structure s of arbitrary species \(\Sigma\) (cf.~Ens., chap.~IV, §1, n° 4) and suppose that s is invariant under G; there then exists on \(\mathbf{F}_{b}\) a structure \(s_{b}\) of species \(\Sigma\) and exactly one such that the \(\theta_{x}: F \to F_{b}\) are isomorphisms; one obtains it by transporting s by means of one of the \(\theta_{x}\) (loc. cit., no. 5).

If \(s\) is a section of \(\mathbf{P}\) above an open set \(\mathbf{U}\) of \(\mathbf{B}\) , the map \((b, f) \mapsto s(b)\) . \(f\) is an isomorphism of \(\mathbf{U} \times \mathbf{F}\) onto \(\pi_{\mathbf{E}}^{-1}(\mathbf{U})\) .

6.5.3. Examples:

\begin{enumerate}
\def\labelenumi{(\alph{enumi})}
\tightlist
\item
  Let \(\lambda = (\mathbf{B} \times \mathbf{G}, \mathbf{B}, \mathrm{pr}_1)\) , let \(\mathbf{E} = \mathbf{B} \times \mathbf{F}\) , and let
\end{enumerate}

\[
\rho : (\mathbf {B} \times \mathbf {G}) \times \mathbf {F} \rightarrow \mathbf {E}
\]

be the map \((b, g, f) \mapsto (b, g \cdot f)\) . We thus obtain on \(\mathbf{B} \times \mathbf{F}\) a fiber bundle structure associated with \(\lambda\) with typical fiber \(\mathbf{F}\) , which is said to be trivial.

\begin{enumerate}
\def\labelenumi{(\alph{enumi})}
\setcounter{enumi}{1}
\item
  Let \(\lambda = (\mathbf{P},\mathbf{G},\mathbf{B},\pi)\) and \(\lambda' = (\mathbf{P}',\mathbf{G}',\mathbf{B}',\pi')\) be two principal fibrations. Let F (resp. F') be a variety on which G (resp. G') acts on the left; and let E (resp. E') be a fiber space associated with \(\lambda\) (respectively to \(\lambda'\) ) with fiber type F (respectively F'). The group \(\mathbf{G} \times \mathbf{G}'\) acts on \(\mathbf{F} \times \mathbf{F}'\) by \((g, g') \cdot (f, f') = (g \cdot f, g' \cdot f')\) . The map \((\mathbf{P} \times \mathbf{P}') \times (\mathbf{F} \times \mathbf{F}') \to \mathbf{E} \times \mathbf{E}'\) is the product of the frame maps \(\mathbf{P} \times \mathbf{F} \to \mathbf{E}\) and \(\mathbf{P}' \times \mathbf{F}' \to \mathbf{E}'\) and endows \(\mathbf{E} \times \mathbf{E}'\) with a structure of fiber bundle associated with \(\lambda \times \lambda'\) with typical fiber \(\mathbf{F} \times \mathbf{F}'\) .
\item
  With the notations of (b), if one supposes that \(B' = B\) , one defines in the same way on \(E \times_{B} E'\) a fiber bundle structure associated with \(\lambda \times_{B} \lambda'\) with typical fiber \(F \times F'\) .
\end{enumerate}

6.5.4. The notations being those of no. 6.5.1, let \(h: \mathbf{B}' \to \mathbf{B}\) be a morphism, let \(\lambda' = \mathbf{B}' \times_{\mathbf{B}} \lambda\) , and let \(\mathbf{E}' = \mathbf{B}' \times_{\mathbf{B}} \mathbf{E}\) . If \(\mathbf{P}' = \mathbf{B}' \times_{\mathbf{B}} \mathbf{P}\) , define a map \(\mathbf{P}' \times \mathbf{F} \to \mathbf{E}'\) by setting \((b', x).f = (b', x.f)\) ; one thus endows \(\mathbf{E}'\) with a structure of fiber bundle associated with \(\lambda'\) with fiber type F; it is called the inverse image by \(h\) of the structure given on E.

6.5.5. Let \(\lambda = (\mathbf{P}, \mathbf{G}, \mathbf{B}, \pi)\) be a principal fibration, and let E (resp. E') be a fiber space associated with \(\lambda\) with fiber type F (resp. F'). Let \(u: \mathbf{F} \to \mathbf{F}'\) be a morphism compatible with the operations of G (that is, such that \(u(g.f) = g \cdot u(f)\) for \(f \in \mathbf{F}\) , \(g \in \mathbf{G}\) ). There then exists a unique morphism \(\bar{u}: \mathbf{E} \to \mathbf{E}'\) such that one has \(\bar{u}(x.f) = x \cdot u(f)\) for \(x \in \mathbf{P}\) , \(f \in \mathbf{F}\) . If \(u\) is an immersion (resp. a submersion, a subimmersion), the same is true of \(\bar{u}\) .

In particular, suppose that a group variety H operates on the right on F in such a way that \(g.(f.h) = (g.f).h\) for \(g \in G\) , \(f \in F\) , \(h \in H\) . Every \(h \in H\) then defines an automorphism \(u_h\) of F compatible with the operations of G, whence an automorphism \(\bar{u}_h\) of E; the group H operates on the right on E by \((y,h) \mapsto \bar{u}_h(y)\) .

6.5.6. Let \(\lambda = (\mathbf{P},\mathbf{G},\mathbf{B},\pi)\) be a principal fibration, and let E be a fiber space associated with \(\lambda\) with fiber type F. Let \(s: \mathbf{B} \to \mathbf{E}\) be a section of E. For every \(x \in \mathbf{P}\) , there exists a unique element \(\sigma(x) \in \mathbf{F}\) such that \(s(\pi(x)) = x \cdot \sigma(x)\) . The map \(\sigma: \mathbf{P} \to \mathbf{F}\) thus defined is a morphism of varieties, and satisfies the identity:

\[
\sigma (x. g) = g ^ {- 1}. \sigma (x).\tag{*}
\]

The map \(s \mapsto \sigma\) is a bijection of the set of sections of E onto the set of morphisms of P into F which satisfy the identity (*).

6.5.7. Let \(\lambda = (\mathbf{P}, \mathbf{G}, \mathbf{B}, \pi)\) and \(\lambda' = (\mathbf{P}', \mathbf{G}, \mathbf{B}, \pi')\) be two principal fibrations with the same base \(\mathbf{B}\) and of the same structural group \(\mathbf{G}\) . The principal fibration \(\lambda \times_{\mathbf{B}} \lambda' = (\mathbf{P} \times_{\mathbf{B}} \mathbf{P}', \mathbf{G} \times \mathbf{G}, \mathbf{B}, (\pi, \pi')_{\mathbf{B}})\) has as its structure group \(\mathbf{G} \times \mathbf{G}\) . Let us make \(\mathbf{G} \times \mathbf{G}\) operate on the left on \(\mathbf{G}\) by the formula:

\[
(g, g ^ {\prime}) \cdot g _ {1} = g \cdot g _ {1} \cdot g ^ {- 1},
\]

and let E be the fiber bundle associated with \(\lambda \times_{\mathbf{B}} \lambda'\) with fiber type G (endowed with the operation law defined above). Then the sections of E correspond bijectively to the isomorphisms from P to P'. Precisely, if s is a section of E corresponding (cf.~no. 6.5.6) to the morphism \(\sigma: \mathbf{P} \times_{\mathbf{B}} \mathbf{P}' \to \mathbf{G}\) , there exists a G-B-isomorphism \(f_s: \mathbf{P} \to \mathbf{P}'\) such that \(\sigma(x, f_s(x)) = e\) for all \(x \in \mathbf{P}\) ; the map \(s \mapsto f_s\) is a bijection from the set of sections of E onto the set of G-B-isomorphisms from P to P'.

\subsection{Extension and restriction of the structure group}

6.6.1. Let \(\lambda = (\mathbf{P}, \mathbf{G}, \mathbf{B}, \pi)\) be a principal fibration, and let \(\varphi\) be a homomorphism from \(\mathbf{G}\) into a group variety \(\mathbf{H}\) . Let us make \(\mathbf{G}\) operate on the left on \(\mathbf{H}\) by \(g \cdot h = \varphi(g) \cdot h\) and let \(\mathbf{P} \times^{\mathbf{G}} \mathbf{H}\) be the associated fiber bundle to \(\lambda\) with typical fiber \(\mathbf{H}\) . Since the right translations in \(\mathbf{H}\) are compatible with the operations of \(\mathbf{G}\) , the group \(\mathbf{H}\) operates on the right on \(\mathbf{P} \times^{\mathbf{G}} \mathbf{H}\) (cf.~no. 6.5.5); if \(\pi_{\mathbf{H}}\) denotes the projection of \(\mathbf{P} \times^{\mathbf{G}} \mathbf{H}\) onto \(\mathbf{B}\) , the quadruple ( \(\mathbf{P} \times^{\mathbf{G}} \mathbf{H}\) , \(\mathbf{H}\) , \(\mathbf{B}\) , \(\pi_{\mathbf{H}}\) ) is a principal fibration, denoted \(\varphi(\lambda)\) ; it is said to be deduced from \(\lambda\) by the homomorphism \(\varphi\) .

The map \(f\) of \(\mathbf{P}\) to \(\mathbf{P} \times^{\mathbf{G}}\mathbf{H}\) which associates to \(x \in \mathbf{P}\) the class of \((x, e)\) is a \(\mathbf{B}\) -morphism of \(\mathbf{P}\) to \(\mathbf{P} \times^{\mathbf{G}}\mathbf{H}\) compatible with \(\varphi\) (cf.~no. 6.3.1). Moreover, if \(f'\) is a \(\mathbf{B}\) -morphism of \(\mathbf{P}\) into a principal fiber bundle \(\mathbf{P}'\) of the structural group \(\mathbf{H}\) and if \(f'\) is compatible with \(\varphi\) , there exists a unique \(\mathbf{H}\) - \(\mathbf{B}\) -isomorphism \(\theta\) of \(\mathbf{P} \times^{\mathbf{G}}\mathbf{H}\) onto \(\mathbf{P}'\) such that \(f' = \theta \circ f\) .

6.6.2. Suppose that \(\lambda\) is defined by means of an open covering \(\mathcal{V} = (\mathrm{U}_i)\) of B and of a cocycle \((g_{ij})\) (6.4.2). Then \(\varphi(\lambda)\) can be defined using the same covering and the cocycle \((h_{ij})\) , with \(h_{ij} = \varphi \circ g_{ij}\) .

6.6.3. Let F be a variety on which the group H acts on the left; we denote by \((h, y) \mapsto h \cdot y\) the action law of H on F. The group G acts on F by \((g, y) \mapsto \varphi(g) \cdot y\) . Let E be a fiber bundle associated with \(\varphi(\lambda)\) with typical fiber F. The map \((x, y) \mapsto f(x) \cdot y\) from P × F to E (the map f being that defined in no. 6.6.1) endows E with a structure of fiber bundle associated with \(\lambda\) of fiber type F. In particular, \((\mathbf{P} \times^{\mathbf{G}}\mathbf{H}) \times^{\mathbf{H}}\mathbf{F}\) is identified with P × \(^{\mathbf{G}}\) F.

6.6.4. Suppose that G is a Lie subgroup of H, and that \(\varphi: \mathbf{G} \to \mathbf{H}\) is the canonical injection of G into H. One then says that \(\varphi(\lambda)\) is deduced from \(\lambda\) by extension to H of the structural group. The morphism \(f: \mathbf{P} \to \mathbf{P} \times^{\mathbf{G}} \mathbf{H}\) of No. 6.6.1 is an isomorphism from P onto a closed submanifold of \(\mathbf{P} \times^{\mathbf{G}} \mathbf{H}\) (equal to \(\mathbf{P} \times^{\mathbf{G}} \mathbf{H}\) if \(\mathbf{G} = \mathbf{H}\) ); this isomorphism is compatible with the operations of G (note that G acts on \(\mathbf{P} \times^{\mathbf{G}} \mathbf{H}\) as a subgroup of H).

6.6.5. Suppose further that G is a Lie subgroup of H, and let \(\mu = (Q, H, B, \pi)\) be a principal fibration with structure group H and base B, and let E be a fiber bundle associated with \(\mu\) of fiber type H/G. If \(\gamma\) denotes the canonical projection of H onto H/G, set \(\delta(y) = y \cdot \gamma(e)\) , for \(y \in Q\) (where \(e\) denotes the identity element of H); one thus obtains a morphism \(\delta: Q \to E\) and the quadruple \((Q, G, E, \delta)\) is a principal fibration. In particular, Q/G is canonically identified with E, itself isomorphic to \(Q \times^{H} H/G^{1}\) .

Now let \(s: \mathbf{B} \to \mathbf{E}\) be a section of \(\mathbf{E}\) , and let \(\lambda_s\) be the inverse image under \(s\) of the fibration \((\mathbf{Q}, \mathbf{G}, \mathbf{E}, \delta)\) that we have just defined. It is a principal fibration, with base \(\mathbf{B}\) and with structural group \(\mathbf{G}\) , and its extension to \(\mathbf{H}\) is isomorphic to \(\mu\) ; any principal fibration enjoying these properties can be obtained in this way, up to isomorphism; two sections \(s_1\) and \(s_2\) of \(\mathbf{E}\) define isomorphic fibrations if and only if they are transformed into one another by an H-B-automorphism of \(\mu\) .

\subsection{Changes of structure}

The structures and operations described in this paragraph are compatible with the changes of structure described in numbers 5.13 and 5.14.

\section{Vector bundles}

Throughout this paragraph, the letter B denotes a manifold of class \(C^{r}\) ( \(r \geqslant 1\) ) and the letter M denotes a set equipped with a map \(\pi\) from M to B. One says that B is the base of M and, for every \(b \in B\) , one denotes the fiber of M at b by \(M_{b}\) and calls it the subset \(\pi^{-1}(b)\) of M.
